\documentclass[11pt,reqno]{amsart}

\usepackage[T1]{fontenc}
\usepackage{lmodern}
\usepackage{microtype}
\usepackage{amsmath,amssymb,mathtools,mathrsfs}
\usepackage{booktabs,longtable,array,tabularx}
\usepackage{enumitem}
\usepackage{xcolor}
\usepackage{listings}
\usepackage{aliascnt}
\usepackage[colorlinks=true,
  linkcolor=blue!55!black,
  citecolor=green!40!black,
  urlcolor=blue!60!black,
  pdftitle={A Prescribed-Factor Sieve and an Unconditional Prime-Gap Bound of 216},
  pdfauthor={Jingwei Liu},
  pdfsubject={Equidistribution to prescribed-factor moduli and bounded gaps between primes},
  pdfkeywords={bounded gaps between primes, Maynard sieve, q-van der Corput method, equidistribution, trace functions, interval arithmetic}]{hyperref}
\usepackage[nameinlink,noabbrev]{cleveref}

\allowdisplaybreaks
\numberwithin{equation}{section}
\setlist{nosep}
\raggedbottom

\newtheorem{theorem}{Theorem}[section]
\newaliascnt{proposition}{theorem}
\newtheorem{proposition}[proposition]{Proposition}
\aliascntresetthe{proposition}
\newaliascnt{lemma}{theorem}
\newtheorem{lemma}[lemma]{Lemma}
\aliascntresetthe{lemma}
\newaliascnt{corollary}{theorem}
\newtheorem{corollary}[corollary]{Corollary}
\aliascntresetthe{corollary}
\newaliascnt{hypothesis}{theorem}
\newtheorem{hypothesis}[hypothesis]{Hypothesis}
\aliascntresetthe{hypothesis}
\theoremstyle{definition}
\newaliascnt{definition}{theorem}
\newtheorem{definition}[definition]{Definition}
\aliascntresetthe{definition}
\newaliascnt{remark}{theorem}
\newtheorem{remark}[remark]{Remark}
\aliascntresetthe{remark}

\crefname{theorem}{Theorem}{Theorems}
\crefname{proposition}{Proposition}{Propositions}
\crefname{lemma}{Lemma}{Lemmas}
\crefname{corollary}{Corollary}{Corollaries}
\crefname{hypothesis}{Hypothesis}{Hypotheses}
\crefname{definition}{Definition}{Definitions}
\crefname{remark}{Remark}{Remarks}
\Crefname{theorem}{Theorem}{Theorems}
\Crefname{proposition}{Proposition}{Propositions}
\Crefname{lemma}{Lemma}{Lemmas}
\Crefname{corollary}{Corollary}{Corollaries}
\Crefname{hypothesis}{Hypothesis}{Hypotheses}
\Crefname{definition}{Definition}{Definitions}
\Crefname{remark}{Remark}{Remarks}

\newcommand{\N}{\mathbb N}
\newcommand{\R}{\mathbb R}
\newcommand{\Pp}{\mathbb P}
\newcommand{\e}{\mathrm e}
\newcommand{\dd}{\,\mathrm d}
\newcommand{\1}{\mathbf 1}
\newcommand{\supp}{\operatorname{supp}}
\newcommand{\lcm}{\operatorname{lcm}}
\newcommand{\rad}{\operatorname{rad}}
\newcommand{\DHL}{\operatorname{DHL}}
\newcommand{\SW}{\operatorname{SW}}
\newcommand{\DS}{\mathcal D_{\mathrm{pf}}}
\newcommand{\Qfam}{\mathcal Q^{\sharp}}
\newcommand{\Rfam}{\mathcal R^{\sharp}}
\newcommand{\eps}{\varepsilon}

\title[An unconditional prime-gap bound of 216]
{A Prescribed-Factor Sieve and an Unconditional\\
Prime-Gap Bound of $216$}

\author{Jingwei Liu}
\date{September 3, 2026}

\subjclass[2020]{Primary 11N05, 11N36; Secondary 11L07, 11N35, 65G30}
\keywords{bounded gaps between primes, Maynard--Tao sieve,
equidistribution in arithmetic progressions, q-van der Corput method,
trace functions, rigorous interval arithmetic}

\begin{document}

\begin{abstract}
Let $H_1=\liminf_{n\to\infty}(p_{n+1}-p_n)$.  We prove the
unconditional bound $H_1\le216$ by combining a prescribed-factor
equidistribution theorem with a rigorously certified Maynard sieve weight.
The analytic theorem treats Type-I convolutions in arithmetic progressions
to squarefree moduli $d\le x^{1/2+2\omega}$ possessing an explicit chain of
divisors in prescribed scale windows.  Its proof replaces global smoothness
of the modulus by transport of these divisor witnesses, retains both branches
of the genuine terminal scale, and closes the supported $q$-van der Corput
descent and outer dispersion argument without a conjectural input.

The finite sieve component uses a layered support region recording both the
number and total mass of coordinates above a fixed threshold.  A
$208$-parameter rational test function for $k=46$ is verified by exact
support arithmetic and $1024$-bit outward-rounded Arb integration, giving
\[
\frac{46J(F)}{I(F)}>
1.0000266555919862183480792294536290427.
\]
We also certify an admissible $46$-tuple of diameter $216$.  The
prescribed-factor theorem, the exact support decomposition, and the
Maynard--Polymath transfer then give infinitely many pairs of primes at
distance at most $216$.  The paper distinguishes the cited trace and
dispersion theorems from the new support-transport argument; no unproved
equidistribution hypothesis is used.  All finite data, hashes, interval
certificates, and verification programs are supplied as ancillary files.
\end{abstract}

\maketitle
\tableofcontents

\section{Introduction}

Let $p_n$ denote the $n$th prime and put
\[
 H_m=\liminf_{n\to\infty}(p_{n+m}-p_n).
\]
Zhang's proof that $H_1$ is finite initiated a rapid sequence of
improvements based on distribution estimates for primes in arithmetic
progressions and multidimensional sieve weights
\cite{Zhang2014,Maynard2015,PolymathEDZ2014,PolymathSieve2014}.
The Polymath project obtained the unconditional bound $H_1\le246$.
Stadlmann subsequently strengthened the smooth-modulus q-van der Corput
estimate \cite{Stadlmann2025}; a recent preprint proposes combining related
equidistribution estimates with a layered sieve support to deduce
$H_1\le240$ \cite{Stadlmann2026}.

The purpose of this paper is twofold.  First, we prove the exact
prescribed-factor Type-I statement needed to enlarge the sieve support.  The
proof starts from the established dispersion and trace-function estimates of
D.~H.~J. Polymath, replaces every use of global modulus smoothness by an
explicit factor or squarefree-support argument, and closes the return from
the terminal descendants to the original discrepancy.  Second, we construct
and rigorously certify a layered $k=46$ Maynard weight.  Together these two
components yield the following result.

\begin{theorem}[Prime-gap bound]\label{thm:main}
One has
\[
 \boxed{H_1\le216.}
\]
Equivalently, infinitely many pairs of distinct primes differ by at most
$216$.
\end{theorem}

The new analytic input is stated next.  The terminology concerning
coefficient sequences and equidistribution is recalled in
\cref{sec:preliminaries}.

\begin{definition}[Prescribed-factor family]\label{def:S52-family}
For fixed $\omega,\delta,\gamma,\eta>0$, let $\DS(x)$ consist of the
integers $d\le x^{1/2+2\omega}$ for which there exist
\[
 r\mid d,\qquad u\mid d/r,\qquad d_1\mid r
\]
satisfying
\begin{align}
 x^{\gamma-3\eta-\delta}<r&<x^{\gamma-3\eta},\label{eq:S52-r}\\
 x^{1-\gamma-6\eta-\delta}d^{-1}<u
 &<x^{1-\gamma-6\eta}d^{-1},\label{eq:S52-u}\\
 r^2x^{2-\gamma-52\eta-\delta}d^{-4}<d_1
 &<r^2x^{2-\gamma-52\eta}d^{-4}.
 \label{eq:S52-d1}
\end{align}
\end{definition}

\begin{theorem}[Prescribed-factor Type-I estimate]
\label{thm:S52}
Let $f=\alpha\star\beta$, where $\alpha$ and $\beta$ are coefficient
sequences at scales $M$ and $N=x^\gamma$, respectively,
$MN\asymp x$, $N\le x^{1/2}$, and $\beta$ has the Siegel--Walfisz
property.  Suppose
\begin{align}
 8\omega+4\delta+2\gamma&<1,\label{eq:S52-a}\\
 32\omega+10\delta-\gamma&<0,\label{eq:S52-b}\\
 48\omega+16\delta-4\gamma&<-1.\label{eq:S52-c}
\end{align}
There is $\eta_0=\eta_0(\omega,\delta,\gamma)>0$ such that, whenever
$0<\eta<\eta_0$, for every $A>0$ and every integer $a=a(x)$ satisfying
$(a,p)=1$ for every prime $p\le x$,
\begin{equation}\label{eq:S52-conclusion}
 \sum_{\substack{d\in\DS(x)\\d\ \mathrm{squarefree}}}
 \left|
 \sum_{\substack{x\le n\le2x\\n\equiv a\pmod d}}f(n;x)
 -\frac1{\varphi(d)}
 \sum_{\substack{x\le n\le2x\\(n,d)=1}}f(n;x)
 \right|
 \ll_{A,\eta,\omega,\delta,\gamma}
 \frac{x}{(\log x)^A}.
\end{equation}
The estimate is uniform in the permitted residue $a$, and the constants are
uniform when $(\omega,\delta,\gamma)$ ranges over a compact subset of the
open region above.
\end{theorem}

Our second input is a finite theorem.

\begin{theorem}[Layered variational certificate]\label{thm:finite-intro}
For $k=46$, $R=521/2000$, $\delta=1/50$, and
$\omega=57/10000$, there exists a smooth symmetric Maynard test function
$F$ supported on the layered region defined in \cref{def:layered-region}
such that
\begin{equation}\label{eq:finite-intro}
 \frac{46J(F)}{I(F)}>
 1.0000266555919862183480792294536290427.
\end{equation}
The strict inequality is certified by exact rational support arithmetic and
outward-rounded $1024$-bit interval integration.
\end{theorem}

\subsection{The correction completed in this paper}

The prescribed-factor statement in \cref{thm:S52} appears as Lemma~6 of
\cite[version~1, p.~12]{Stadlmann2026}.  The published 2025 theorem
\cite{Stadlmann2025}, however, applies to globally $x^\delta$-smooth moduli.
Passing to the larger prescribed-factor family requires an additional
inheritance argument.

Two concrete issues prevent the extension sketch in
\cite[version~1, pp.~12, 17--18]{Stadlmann2026} from supplying that argument
as written.  First, the modulus family is stated with a $52\eta$ window,
restated in the proof with a $100\eta$ window, and then used again with
$52\eta$.  Second, the proof introduces
\[
 \Delta^*=\min\left\{\frac{N}{|\Lambda|x^{5\eta}},\Delta_1\right\}
\]
and sets $\Delta^*=\Delta_1$ using two quantities that are both bounded
below by a third.  Such bounds do not compare the two quantities and hence do
not justify the asserted identity.  We discuss these points precisely in
\cref{sec:diagnosis}.  Our repair keeps the $52\eta$ family and retains both
branches of $\Delta^*$ throughout.

No conclusion is drawn from criticism alone.  The terminal trace estimate,
support transport, supported descent, two-branch exponent ledger, and outer
dispersion closure are proved in
\cref{sec:terminal,sec:interfaces,sec:descent,sec:closure,sec:outer}.  Their
composition proves \cref{thm:S52} without using the prescribed-factor claim
of \cite{Stadlmann2026} as an input.

\subsection{Main novelties}

The paper introduces four ingredients.
\begin{enumerate}[label=\textup{(\roman*)}]
\item A long-scale completion lemma for the terminal two-variable trace sum,
uniform for arbitrary squarefree polynomial-size moduli.  It uses the second
estimate in Polymath Proposition~8.4 and does not require dense divisibility.
\item A factor-interface calculus that replaces the two uses of global
smoothness by explicit divisors in prescribed windows.  Prime supports,
units, valuation saturation, $q_0$-powers, and subpower preprocessing losses
are transported separately.
\item An explicit one-way inheritance proof.  It fixes the order of witness
selection, Cauchy duplication, nonnegative enlargement, terminal completion,
and outer Cartesian summation.  This closes the local and outer transitions
without transferring a witness to an unrelated Cauchy copy.
\item A layered variational ansatz that records both the count and mass of
coordinates above $\delta$.  This retains information lost by a purely
radial ansatz and produces a certified crossing with $k=46$ on the corrected
support.
\end{enumerate}

\subsection{Computer assistance and logical scope}

The analytic proof is written independently of the verification programs.
The programs check finite exponent identities, support incidence relations,
and parameter margins; they do not replace the cited trace estimates or the
analytic arguments proving \cref{thm:S52}.  Computer assistance is essential
for \cref{thm:finite-intro}.  There every input is rational, the support cover
is checked by exact linear arithmetic, and all integrals are enclosed by Arb
intervals.  The candidate and certificates are linked by SHA-256 hashes.

\subsection{Organization}

After the preliminaries, \cref{sec:diagnosis} compares the smooth-modulus
parent theorem with the prescribed-factor extension and identifies the two
transitions requiring repair.  Sections \ref{sec:terminal}--\ref{sec:outer}
prove \cref{thm:S52}.  Section \ref{sec:layered} proves the finite certificate,
and \cref{sec:main-proof} combines the analytic and finite components to prove
\cref{thm:main}.  The appendices record the dependency graph and the
reproducibility manifest.

\section{Preliminaries}\label{sec:preliminaries}

Throughout, $x$ tends to infinity.  Implied constants may depend on fixed
parameters and fixed seminorm bounds, but not on $x$ or on the particular
coefficient sequences under consideration.  We write
$e_q(z)=\exp(2\pi iz/q)$ and use $\tau$ for the divisor function.

\begin{definition}[Coefficient sequences]\label{def:coefficients}
A coefficient sequence is a function
$\alpha:\N\times(1,\infty)\to\R$ satisfying
\[
 |\alpha(n;x)|\ll \tau(n)^{O(1)}(\log x)^{O(1)}.
\]
It is \emph{located at scale} $N=N(x)$ if, for fixed $0<c<C$, it is
supported on $[cN,CN]$.  A sequence at scale $N$ has the
\emph{Siegel--Walfisz property} if, for all $A>0$, all $q,r\ge1$, and
all $(a,q)=1$,
\begin{equation}\label{eq:sw}
 \left|
 \sum_{\substack{n\equiv a\pmod q\\(n,r)=1}}\alpha(n;x)
 -\frac1{\varphi(q)}\sum_{(n,qr)=1}\alpha(n;x)
 \right|
 \ll_A \tau(qr)^{O(1)}\frac{N}{(\log x)^A}.
\end{equation}
It is \emph{smooth at scale} $N$ if
$\alpha(n;x)=\psi_x(n/N)$, where all $\psi_x$ are supported in one fixed
compact interval and every fixed derivative is bounded by a fixed power of
$\log x$.
\end{definition}

For $f=\alpha\star\beta$ put
\begin{equation}\label{eq:discrepancy}
 \Delta_f(a;d)=
 \sum_{\substack{x\le n\le2x\\n\equiv a\pmod d}}f(n;x)
 -\frac1{\varphi(d)}
 \sum_{\substack{x\le n\le2x\\(n,d)=1}}f(n;x).
\end{equation}
Thus \cref{eq:S52-conclusion} is the assertion
\[
 \sum_{\substack{d\in\DS\\d\ \mathrm{squarefree}}}
 |\Delta_f(a;d)|\ll_Ax(\log x)^{-A}.
\]

\subsection{Established analytic inputs}

We use two groups of previously established results.

\begin{enumerate}[label=\textup{(\roman*)}]
\item\label{input:trace} The sheaf-theoretic complete-sum and Fourier-transform estimates in
\cite[Theorem~6.17, Corollary~6.24, and Proposition~8.4]{PolymathEDZ2014}.
Only the second bound in Proposition~8.4 is used at the terminal stage.  Its
local rational phase is nondegenerate whenever the two denominator slopes are
distinct modulo every prime factor of the squarefree modulus.
\item\label{input:outer} The dispersion identity and initial Type-I reduction
in \cite[Sections~5 and 8A]{PolymathEDZ2014}, in the explicit form reorganized
for smooth moduli in \cite{Stadlmann2025}.  We use the identity before any
smoothness-based enlargement and prove the required enlargement for $\DS$.
\item\label{input:qvdc} The algebraic q-van der Corput identities from
\cite{Stadlmann2025}: Poisson completion, Cauchy--Schwarz, finite M\"obius
inversion, CRT substitutions, and Taylor expansion of fixed smooth profiles.
Every implication needed for the prescribed-factor family is stated and
proved below; the smooth-modulus conclusion itself is not used as an input.
\end{enumerate}

The Bombieri--Vinogradov theorem, the prime number theorem, and the standard
Maynard--Polymath multidimensional Selberg sieve are used in their customary
forms \cite{Maynard2015,PolymathSieve2014}.

\subsection{Parameter reserve}

Assume \cref{eq:S52-a,eq:S52-b,eq:S52-c}.  Set
\begin{align}
\mathfrak m=\min\{&1-8\omega-4\delta-2\gamma,
 \gamma-32\omega-10\delta,
 4\gamma-1-48\omega-16\delta,\notag\\
&\tfrac12-2\omega-\gamma,
 \gamma-\delta,
 \gamma-8\omega-4\delta,
 \tfrac14-16\omega-6\delta\}.
\label{eq:margin}
\end{align}
Every entry is positive.  We fix
\begin{equation}\label{eq:eta-choice}
 0<\eta<\eta_0:=
 \min\left\{\frac{\delta}{10^{100}},\frac{\mathfrak m}{1600}\right\},
 \qquad \eps_{\rm tr}=\frac{\eta}{1000}.
\end{equation}
The parameter $\eta$ pays structural losses in the q-van der Corput chain;
$\eps_{\rm tr}$ is reserved for conductor, divisor, and finite-seminorm
losses in the terminal trace estimate.  Keeping them separate prevents the
same power saving from being spent twice.

The first and third inequalities in the parameter region imply
\begin{equation}\label{eq:aux-range}
 12\omega+6\delta<\gamma<\frac12-2\omega.
\end{equation}
All later auxiliary range conditions follow from
\cref{eq:S52-a,eq:S52-b,eq:S52-c,eq:eta-choice}.

\section{The prescribed-factor extension and the point requiring repair}
\label{sec:diagnosis}

This section records the precise relation between the published
smooth-modulus theorem \cite{Stadlmann2025}, the S.~Type-I extension proposed
in \cite{Stadlmann2026}, and \cref{thm:S52}.  The distinction matters
because the larger modulus family is not covered by the published theorem.

\subsection{What the published theorem supplies}

The 2025 theorem assumes that the relevant modulus divides $P(x^\delta)$;
in particular it is squarefree and every prime factor is at most $x^\delta$.
Its q-van der Corput argument introduces a divisor at the scale
\[
 D=\frac{N}{x^{50\eta}H^2}.
\]
This theorem and its local reduction are established analytic inputs in the
present paper.  What must be proved afresh is that the reduction survives when
global smoothness is replaced by the two divisor witnesses in the definition
of $\DS$.

\subsection{The window inconsistency}

In Lemma~6 on page~12 of version~1 of \cite{Stadlmann2026}, the
S.~Type-I statement uses
\[
 r^2x^{2-\gamma-52\eta-\delta}d^{-4}<d_1<
 r^2x^{2-\gamma-52\eta}d^{-4}.
\]
At the beginning of its proof on page~17 this is restated with $100\eta$,
while the factor family immediately following the restatement again uses
$52\eta$.
These three displayed conditions define different modulus families.  The
proof therefore does not, as written, establish the stated lemma.

The compatible exponent is $52$.  Indeed, after the first factorization,
\[
 H=\frac{x^\eta RQ^2}{q_0M},
\]
and the stated divisor window gives
\[
 \frac{x^{2-\gamma-52\eta}}{Q^4R^2}
 \asymp \frac{N}{q_0^2x^{50\eta}H^2}.
\]
Thus $52=50+2$, with the additional $2\eta$ arising from the square of the
$x^\eta$ in the definition of $H$.  We retain this window throughout.

\subsection{The terminal minimum}

The more substantive issue occurs on page~18 and concerns
\begin{equation}\label{eq:delta-star-diagnosis}
 \Delta^*=\min\left\{
 \frac{N}{|\Lambda|x^{5\eta}},\Delta_1\right\}.
\end{equation}
The extension sketch gives lower bounds of the form
\[
 \frac{N}{|\Lambda|x^{5\eta}}\gg C,
 \qquad \Delta_1\gg C,
\]
and then sets $\Delta^*=\Delta_1$.  This conclusion does not follow: two
quantities being bounded below by the same third quantity gives no ordering
between them.  In particular, it does not establish
$\Delta_1\le N/(|\Lambda|x^{5\eta})$.

This is not merely a typographical matter, because reciprocal powers of
$\Delta^*$ occur in the final estimates.  Our repair is to preserve the
minimum and use the exact identity
\begin{equation}\label{eq:delta-branch-diagnosis}
 \frac1{\Delta_1\Delta^*}
 =\max\left\{
 \frac{|\Lambda|x^{5\eta}}{N\Delta_1},
 \frac1{\Delta_1^2}
 \right\}.
\end{equation}
Both branches are carried through the exponent ledger in
\cref{sec:closure}.

\subsection{Scholarly claim boundary}

The observations above do not invalidate the published smooth-modulus theorem
\cite{Stadlmann2025}; they concern only the larger prescribed-factor
extension in \cite{Stadlmann2026}.  Nor do they imply that the desired
extension is false.  Sections \ref{sec:terminal}--\ref{sec:outer} give a
complete replacement argument: the supported local descent is proved first,
and only then consumed by the outer Cartesian and dispersion closure.  The
attribution used below is therefore: the smooth-modulus q-van der Corput
architecture is due to Stadlmann, while the factor-interface transport to
$\DS$, the genuine-minimum branch analysis, the supported descent, and the
corrected outer Cartesian closure are contributions of the present paper.

\section{The terminal squarefree trace estimate}
\label{sec:terminal}

Fix $C\ge1$.  Let $\mathcal P$ be a class of profiles supported in one
fixed compact interval and satisfying

\begin{equation}\label{eq:profile-bounds}
 \|\psi^{(j)}\|_{L^1}+\|\psi^{(j)}\|_{L^\infty}
 \le B_j(\log x)^{A_j},\qquad 0\le j\le2,
\end{equation}

for fixed $A_j,B_j$.  Write
$\psi_T(n)=\psi((n-x_0)/T)$, allowing an arbitrary real shift $x_0$.
When $T<1$, this notation denotes a sequence supported on $O_{\mathcal P}(1)$
integers and bounded by $(\log x)^{O_{\mathcal P}(1)}$.

\begin{lemma}[Uniform long-scale completion]\label{lem:long-completion}
Let $q\le x^C$ be squarefree, and let $K(\cdot;q)$ be a composite trace
function whose prime-local sheaves have bounded conductor and no polynomial
phase of degree at most one.  Uniformly for $1\le T\le x^C$,

\begin{equation}\label{eq:long-completion}
 \left|\sum_n\psi_T(n)K(n;q)\right|
 \ll_{\eps_{\rm tr},\mathcal P}x^{\eps_{\rm tr}}
 \left(q^{1/2}+\frac{T}{q^{1/2}}\right).
\end{equation}
\end{lemma}

\begin{proof}
Periodize $\psi_T$ modulo $q$ and use the normalized finite Fourier
transform

\[
 \widehat K_q(h)=q^{-1/2}\sum_{a\bmod q}K(a;q)e_q(ha).
\]

The Fourier-sheaf estimate in
\cite[Corollary~6.24]{PolymathEDZ2014} gives
$|\widehat K_q(h)|\ll C_0^{\omega(q)}$ for every $h$, including zero.
Squarefreeness and $q\le x^C$ absorb this factor into
$x^{\eps_{\rm tr}/16}$.

Poisson summation and two integrations by parts give, uniformly for
$|y|\le1/2$,

\[
 \left|\sum_n\psi_T(n)e(-ny)\right|
 \ll(\log x)^{O(1)}\frac{T}{(1+T|y|)^2}.
\]

Fourier inversion is exact:

\[
 \sum_n\psi_T(n)K(n;q)
 =q^{-1/2}\sum_{h\bmod q}\widehat K_q(h)
 \sum_n\psi_T(n)e(-hn/q).
\]

The zero frequency is
$O(x^{\eps_{\rm tr}/4}Tq^{-1/2})$, and summation over the nonzero
frequencies is $O(x^{\eps_{\rm tr}/4}q^{1/2})$.  The fixed logarithmic
factors fit in the remaining reserve.
\end{proof}

At a vanishing denominator we use either zero extension or the projective
convention under which a local factor with trivial numerator character is
one.  The following bound is valid under both conventions.

\begin{lemma}[Two-variable squarefree trace bound]
\label{lem:squarefree-trace}
Let $m\le x^C$ be squarefree and let $0<N,\Delta\le x^C$.  For arbitrary
residue classes $A,B,L,\gamma_1,\gamma_2\pmod m$,

\begin{align}
&\left|\sum_d\sum_n\psi_\Delta(d)\psi_N(n)
e_m\!\left(
\frac{AL}{(n+Bd+\gamma_1)(n+(B+L)d+\gamma_2)}
\right)\right|\notag\\
&\quad\ll_{\eps_{\rm tr},\mathcal P}
x^{\eps_{\rm tr}}(AL,m)
\left(\frac N{\sqrt m}+\sqrt m\right)
\left(\frac\Delta{\sqrt m}+\sqrt m\right).
\label{eq:squarefree-trace}
\end{align}
\end{lemma}

\begin{proof}
Assume first that $(AL,m)=1$.  Complete in $n$ modulo $m$.  At a prime
$p\mid m$, CRT gives the local parameter tuple

\[
 (a_p,b_p,c_p,d_p,e_{h,p})
 =\left(B,\gamma_1,B+L,\gamma_2,h(AL)^{-1}\right)\pmod p.
\]

The two slopes differ by $-L\not\equiv0\pmod p$, the numerator character
is nontrivial, and the Fourier parameter is unrestricted.  The local
nondegeneracy hypotheses of
\cite[Theorem~6.17]{PolymathEDZ2014} are therefore satisfied.  Apply
\cref{lem:long-completion} to the completed $n$-sum and then to the
$d$-sum.  This is precisely the squarefree, second estimate in
\cite[Proposition~8.4]{PolymathEDZ2014}.

For $g=(AL,m)$, write $m=gm'$; then $(g,m')=1$ and
$(AL/g,m')=1$.  If $D$ is the product of the two denominators, zero
extension gives

\[
 \mathfrak e_m^{(0)}(AL;D)
 =\mathfrak e_{m'}^{(0)}(AL/g;D)
 \sum_{c\mid g}\mu(c)\1_{c\mid D},
\]

whereas the projective convention simply deletes the trivial local factors.
For squarefree $c$, the condition $c\mid D$ occupies at most
$2^{\omega(c)}c$ pairs of residue classes modulo $c$.  Substitute
$n=cn'+n_0$ and $d=cd'+d_0$ in each pair and apply the coprime estimate
with scales $N/c,\Delta/c$ and modulus $m'$.  If a reduced scale is below
one, direct counting gives the same four-term majorant.  Summing over
$c\mid g$ costs $3^{\omega(g)}=x^{o(1)}$ and yields

\[
 \frac{N\Delta}{m'}+gN+g\Delta+g^2m'
 \ll g\left(\frac N{\sqrt m}+\sqrt m\right)
 \left(\frac\Delta{\sqrt m}+\sqrt m\right).
\]

This proves the result under either convention.
\end{proof}

\begin{lemma}[Congruence-restricted trace bound]
\label{lem:congruence-trace}
Under the hypotheses of \cref{lem:squarefree-trace}, let $q\mid m$ and
restrict $d,n$ to fixed residue classes modulo $q$.  Then

\begin{align}
&\left|\sum_{d\equiv d_*\ (q)}\sum_{n\equiv n_*\ (q)}
\psi_\Delta(d)\psi_N(n)
e_m\!\left(\frac{AL}{(n+Bd)(n+(B+L)d)}\right)\right|\notag\\
&\quad\ll_{\eps_{\rm tr},\mathcal P}
\frac{x^{\eps_{\rm tr}}(AL,m)}q
\left(\frac N{\sqrt m}+\sqrt m\right)
\left(\frac\Delta{\sqrt m}+\sqrt m\right).
\label{eq:congruence-trace}
\end{align}
\end{lemma}

\begin{proof}
Write $d=qd'+d_*$, $n=qn'+n_*$, and $m=qm_1$.  Since $m$ is
squarefree, $(q,m_1)=1$.  The $q$-component is constant; modulo $m_1$
the phase has the same form, with its numerator multiplied by the unit
$\bar q^3$.  Apply \cref{lem:squarefree-trace} at scales $N/q,\Delta/q$
and modulus $m/q$.  The identity

\[
 \frac{N/q}{\sqrt{m/q}}+\sqrt{m/q}
 =q^{-1/2}\left(\frac N{\sqrt m}+\sqrt m\right)
\]

in both variables produces the factor $1/q$.  Direct counting covers a
subunit reduced scale.
\end{proof}


\section{Prescribed factors and prime-support transport}
\label{sec:interfaces}

After the initial Type-I reduction write a contributing modulus as
$d=qr$, where $q\asymp Q$, $r\asymp R$, and put

\begin{equation}\label{eq:H}
 H=\frac{x^\eta RQ^2}{q_0M}.
\end{equation}

The witness $u\mid d/r$, after its common part with $q_0$ is removed,
gives $q_1=u_1v_1$ and the ranges

\begin{align}
 q_0^{-2}x^{-\delta-5\eta}\frac QH\ll U
 &\ll q_0^{-1}x^{-5\eta}\frac QH,\notag\\
 x^{5\eta}H\ll V&\ll q_0x^{\delta+5\eta}H,
 \qquad UV\asymp Q/q_0.
\label{eq:UV}
\end{align}

The second witness and $d\asymp QR$ give

\[
 d_1\asymp\frac{x^{2-\gamma-52\eta}}{Q^4R^2}
 \asymp\frac{N}{q_0^2x^{50\eta}H^2}.
\]

After the standard subdivision into intervals shorter by $x^{5\eta}$,

\begin{equation}\label{eq:Delta1}
 \frac{N}{q_0^2x^{\delta+55\eta}H^2}
 \ll\Delta_1\ll
 \frac{N}{q_0^2x^{55\eta}H^2}.
\end{equation}

\subsection{Reversible small-prime preprocessing}

Put

\[
 D_0=\exp((\log x)^{1/3}),\qquad
 S_x=\exp((\log x)^{2/3}).
\]

\begin{lemma}[Preprocessing]\label{lem:preprocessing}
Let $d$ be squarefree and let $r_0,u_0,d_{1,0}$ witness membership in
$\DS(x)$.  Define

\[
 s(d)=\prod_{\substack{p\mid d\\p\le D_0}}p,\qquad
 q_{\rm raw}=d/r_0,\qquad t=(q_{\rm raw},s(d)).
\]

If $s(d)\le S_x$, then

\[
 q=q_{\rm raw}/t,\qquad r=r_0t
\]

gives $d=qr$ and $(q,P(D_0))=1$.  There exist $u'\mid q$ and
$d_{1,0}\mid r$ satisfying, after dyadic localization,

\begin{align}
 S_x^{-1}\frac{x^{1-\gamma-6\eta-\delta}}{QR}\ll u'
 &\ll\frac{x^{1-\gamma-6\eta}}{QR},\label{eq:pre-u}\\
 S_x^{-2}\frac{x^{2-\gamma-52\eta-\delta}}{Q^4R^2}\ll d_{1,0}
 &\ll\frac{x^{2-\gamma-52\eta}}{Q^4R^2}.
\label{eq:pre-d1}
\end{align}
\end{lemma}

\begin{proof}
Squarefreeness makes the displayed factorization exact.  Every small prime
formerly in $q_{\rm raw}$ divides $t$, so the new $q$ is free of primes at
most $D_0$.  Set $u'=u_0/(u_0,t)$.  Division by at most $S_x$ proves
\cref{eq:pre-u}.  Since $r_0=r/t$, insertion into
\cref{eq:S52-d1} produces exactly the factor $t^{-2}$ and proves
\cref{eq:pre-d1}.
\end{proof}

The exceptional set $s(d)>S_x$ is negligible by the standard
small-prime-tail estimate in the alternative Type-I reduction.  Restoring
the preprocessing changes \cref{eq:UV,eq:Delta1} only by
$S_x^{O(1)}=x^{o(1)}$.  Every later expression has fixed algebraic depth,
so this loss never iterates with $x$.

Define the independent families

\begin{equation}\label{eq:families}
\begin{aligned}
\Qfam(Q,R)=\{q\in[Q,2Q]:{}&(q,P(D_0))=1,\\[-1mm]
&q\text{ has a divisor in \cref{eq:pre-u}}\},\\
\Rfam(Q,R)=\{r\in[R,2R]:{}&
r\text{ has a divisor in \cref{eq:pre-d1}}\}.
\end{aligned}
\end{equation}

Each nonexceptional modulus lies in one selected pair from
$\Qfam\times\Rfam$.  A witness is selected before any Cartesian
enlargement or Cauchy copy is made.

\begin{lemma}[First-factor extraction]\label{lem:first-factor}
Let $q^\sharp=q_0q_1$ be squarefree and let $u\mid q^\sharp$.  If

\[
 a_0=(u,q_0),\qquad u_1=u/a_0,\qquad v_1=q_1/u_1,
\]

then $q_1=u_1v_1$ and $u/q_0\le u_1\le u$.  Consequently the ranges
\cref{eq:UV} hold.
\end{lemma}

\begin{proof}
Prime supports in a squarefree number are sets.  Removing from $u$ the
support shared with $q_0$ leaves a divisor of $q_1$; $v_1$ is its unique
complement.  The inequalities follow from $a_0\le q_0$.
\end{proof}

\begin{lemma}[Second-factor extraction]\label{lem:second-factor}
Choose one witness $d_1\mid r$, write $r=d_1r_1$, and localize
$d_1\asymp\Delta$.  Then

\[
 S_x^{-2}\frac{N}{q_0^2x^{\delta+50\eta}H^2}\ll\Delta
 \ll\frac{N}{q_0^2x^{50\eta}H^2}.
\]

An $O(x^{5\eta})$ smooth cover at scale
$\Delta_1=\Delta/x^{5\eta}$ gives \cref{eq:Delta1}; moreover $d_1,r_1$
are squarefree and coprime.
\end{lemma}

\begin{lemma}[Prime-support transport]\label{lem:prime-support}
Suppose $q_0u_1v_ir$ and $q_0q_2r$ are squarefree for $i=1,2$, and
$(u_1v_1v_2,q_0q_2)=1$.  Let

\[
 r=d_1r_1,\qquad g=(v_1,v_2),\qquad v_i=gv_i^*.
\]

Then the supports of

\[
 d_1,r_1,q_0,u_1,g,v_1^*,v_2^*,q_2
\]

are pairwise disjoint.  Hence

\begin{equation}\label{eq:m}
 m=r_1q_0u_1[v_1,v_2]q_2
 =r_1q_0u_1gv_1^*v_2^*q_2
\end{equation}

is squarefree, $(d_1,m)=1$, and $q_0q_3$ is squarefree for every
$q_3\mid m/q_0$.
\end{lemma}

\begin{proof}
For squarefree integers, divisibility is inclusion of prime supports, gcd is
intersection, lcm is union, and coprimality is disjointness.  The hypotheses
separate all listed supports except the permitted overlap of $v_1,v_2$;
extracting $g$ partitions that last overlap.  Splitting $r=d_1r_1$ proves
the remaining assertions.
\end{proof}

\begin{lemma}[Valuation saturation]\label{lem:valuation}
Let

\[
 m^*=m^{\lfloor2\log x\rfloor},\quad
 w_2=(\lambda,m^*)=(\widetilde\lambda,m^*),\quad
 \lambda^*=(\lambda,\widetilde\lambda).
\]

If $|\lambda|,|\widetilde\lambda|<x$, then, for large $x$,

\[
 (\lambda/w_2,m)=(\widetilde\lambda/w_2,m)
 =(\lambda/\lambda^*,m)=(\widetilde\lambda/\lambda^*,m)=1.
\]
\end{lemma}

\begin{proof}
For every $p\mid m$, one has
$p^{\lfloor2\log x\rfloor}>x$.  Thus the gcd with $m^*$ captures the full
$p$-adic part of each integer.  Removing $w_2$, or the larger common factor
$\lambda^*$, leaves a unit modulo $m$.
\end{proof}

\begin{lemma}[Gcd average without a range hypothesis]
\label{lem:gcd-average}
For $m\ge1$, $K,T>0$, and integers $A\ne0,B$,

\begin{equation}\label{eq:gcd-average}
 \sum_{\substack{|k|\le K\\(Ak+B,m)\le T}}(Ak+B,m)
 \le\tau(m)\bigl(2(A,m)K+T\bigr).
\end{equation}
\end{lemma}

\begin{proof}
Use $(u,m)\le\sum_{c\mid(u,m)}c$, interchange the two finite sums, and
enlarge equality of the gcd to $c\mid Ak+B$.  If soluble, this congruence
is one class modulo $c/(A,c)$ and contributes at most
$2K(A,c)/c+1$.  Multiplication by $c$ gives at most
$2(A,m)K+T$ for every $c\le T$.  There are at most $\tau(m)$ such divisors.
\end{proof}

\section{The supported local descent}
\label{sec:descent}

This section expands the middle of the argument.  In particular, none of the
implications below is justified by saying that the proof for smooth moduli is
unchanged.  We retain the support needed at the next stage, state every
enlargement in the direction in which it is used, and record its loss.

Fix one nonempty outer localization.  Thus $Q,R,q_0,H,H^*,U,V$ satisfy
\cref{eq:H,eq:UV,eq:Delta1}, with $1\le H^*\ll H$, and put

\begin{equation}\label{eq:g-g0}
 g=(v_1,v_2),\qquad g_0=(q_0,\ell),\qquad
 v_i=gv_i^*,\qquad (v_1^*,v_2^*)=1.
\end{equation}

All coefficient sequences in this section are bounded by
$(\log x)^{O(1)}$.  Such factors, as well as a fixed number of dyadic
partitions, are written as $x^{o(1)}$.  Since $\eta$ is fixed, each of them
is smaller than $x^{\eta/100}$ for sufficiently large $x$.

\subsection{Preparation and the first phase identity}

Choose the witness $d_1\mid r$ before duplicating any variable, write
$r=d_1r_1$, and localize $d_1\asymp\Delta$.  In the notation of
\cref{lem:second-factor},

\begin{equation}\label{eq:Delta-range-descent}
 S_x^{-2}\frac{N}{q_0^2x^{\delta+50\eta}H^2}
 \ll\Delta\ll
 \frac{N}{q_0^2x^{50\eta}H^2},
 \qquad \Delta_1=\Delta x^{-5\eta}.
\end{equation}

For fixed $r_1,u_1,v_1,v_2,q_2,d_0$, let $\Sigma_2$ denote

\begin{align}
\Sigma_2={}&
 \sum_{\substack{d\\(d,m)=1\\d\ {\rm squarefree}}}
 \psi_{\Delta_1}(d-d_0)
 \Bigg|
 \sum_{\substack{h_1,h_2\sim H^*\\h_1v_2\ne h_2v_1}}
 \sum_{n}^{\dagger}
 C(n;d)\psi_N(n)\notag\\[-1mm]
 &\hspace{23mm}\times
 \varphi_1(h_1)\overline{\varphi_2(h_2)}
 \Psi(n,h_1,h_2;d)
 \Bigg|,
 \label{eq:Sigma2}
\end{align}

where

\begin{equation}\label{eq:dagger2}
 m=r_1q_0u_1[v_1,v_2]q_2,\qquad
 (n,dr_1q_0u_1v_1v_2)=1,\qquad
 (n+\ell dr_1,q_0q_2)=1.
\end{equation}

The indicator $C(n;d)$ is supported on at most $g_0$ classes of $n$ modulo
$q_0$ for each class of $d$ modulo $q_0$.

\begin{lemma}[Prepared-factor implication]\label{lem:prepared-factor}
Uniformly in the fixed outer data, the estimate

\begin{equation}\label{eq:Sigma2-target}
 \Sigma_2\ll q_0g_0\Delta Ngx^{-10\eta}
\end{equation}

implies

\begin{equation}\label{eq:Sigma1-target}
 \Sigma_1^\sharp\ll g_0RQNUV^2x^{-4\eta}.
\end{equation}
\end{lemma}

\begin{proof}
Grouping the original $r$-sum according to $r=d_1r_1$, followed by dyadic
localization, gives

\begin{equation}\label{eq:prep-group}
 \Sigma_1^\sharp
 \ll x^{\eta/2}\frac{RUQ}{q_0\Delta}
 \sup_{\Delta,r_1,u_1,q_2}\Upsilon_0.
\end{equation}

This is a counting step: the number of choices of the frozen variables is
$O(x^{\eta/2}RUQ/(q_0\Delta))$ after the usual divisor-weight
enlargement.  No factor of a prescribed size is chosen here; the factor was
already supplied by $d_1(r)$.

Separate the relation $h_1v_2=h_2v_1$.  Since
$H^*\ll H\le x^{-5\eta}V$ by \cref{eq:UV}, the number of solutions to that
relation is $O(x^{o(1)}H^*V)$, and its contribution is

\begin{equation}\label{eq:first-diagonal}
 \Upsilon_{0,=}
 \ll x^{o(1)}\Delta H^*VN
 \ll x^{-5\eta+o(1)}\Delta NV^2.
\end{equation}

For the complementary part, cover $d\asymp\Delta$ by
$O(x^{5\eta})$ shifted smooth intervals of length $\Delta_1$.  After the
outer absolute value has been taken, deleting the original family-membership
indicators and the fixed coefficient of modulus at most one can only increase
the sum.  The remaining squarefree and coprimality conditions are exactly
those in \cref{eq:Sigma2}.  Thus \cref{eq:Sigma2-target} and

\begin{equation}\label{eq:gcd-v-average}
 \sum_{v_1\asymp V}\sum_{v_2\asymp V}(v_1,v_2)
 \ll V^2(\log V)^{O(1)}
\end{equation}

give

\[
 \Upsilon_{0,\ne}
 \ll q_0g_0\Delta NV^2x^{-5\eta+o(1)}.
\]

The right side of \cref{eq:first-diagonal} is smaller than this admissible
majorant.  Insert the result in \cref{eq:prep-group}; the factors $q_0$ and
$\Delta$ cancel, and the powers $x^{\eta/2}$ and
$x^{-5\eta+o(1)}$ leave at least $x^{-4\eta}$.  This proves
\cref{eq:Sigma1-target}.
\end{proof}

\begin{lemma}[CRT phase factorization]\label{lem:phase-factorization}
Under the support conditions in \cref{lem:prime-support}, there are residue
classes $A,B\pmod m$, independent of $n,d,h_1,h_2$, such that $(A,m)=1$
and, with $y=h_1v_2^*-h_2v_1^*$,

\begin{equation}\label{eq:phase-factorization}
 \Psi(n,h_1,h_2;d)
 =e_d\!\left(
   \frac{ay}{nr_1q_0u_1gv_1^*v_2^*q_2}
  \right)
  e_m\!\left(\frac{Ay}{d(n+Bd)}\right).
\end{equation}

One may choose $A,B$ by

\begin{align}
 A&\equiv a\pmod{r_1},& B&\equiv0\pmod{r_1},\notag\\
 A&\equiv b_1\pmod{q_0u_1gv_1^*v_2^*},&
 B&\equiv0\pmod{q_0u_1gv_1^*v_2^*},\notag\\
 A&\equiv b_2\pmod{q_2},& B&\equiv\ell r_1\pmod{q_2}.
 \label{eq:AB-CRT}
\end{align}
\end{lemma}

\begin{proof}
By \cref{lem:prime-support}, the three moduli in
\cref{eq:AB-CRT} are pairwise coprime and their product is $m$.
The Chinese remainder theorem therefore defines $A$ and $B$ uniquely modulo
$m$.  Splitting each additive character of the duplicated phase over the
prime-support decomposition and collecting the $d$-component gives the first
factor in \cref{eq:phase-factorization}; collecting the remaining components
gives the second.  The residue $A$ is a unit because $a,b_1,b_2$ are units on
their respective support factors.  This proves the identity and every
inverse appearing in it is legitimate.
\end{proof}

\subsection{The algebraic \texorpdfstring{$q$}{q}-van der Corput step}

Put

\begin{equation}\label{eq:Yset}
 \mathscr Y=\{h_1v_2^*-h_2v_1^*:h_1,h_2\sim H^*\}.
\end{equation}

For a nonzero $y\in\mathscr Y$, set

\[
 w_1=(y,d),\qquad w_2=(y,m^*),\qquad
 m^*=m^{\lfloor2\log x\rfloor}.
\]

After dyadic localization write $w_1\asymp W_1$ and $y\sim Y$, where

\begin{equation}\label{eq:WY-ranges}
 1\ll W_1\ll\Delta,\qquad
 1\ll |Y|\ll\frac{H^*V}{g}.
\end{equation}

For $w_0\mid w_1$, define

\begin{equation}\label{eq:K3-def}
 \mathscr K_3(n,\widetilde n;y,\widetilde y,d)
 :=e_m\!\left(
 \frac{A\{y(\widetilde n+Bd)-\widetilde y(n+Bd)\}}
 {d(n+Bd)(\widetilde n+Bd)}
 \right).
\end{equation}

\begin{align}
\Sigma_3:={}&
 \sum_{\substack{y,\widetilde y\in\mathscr Y\\
 y,\widetilde y\sim Y\\w_1\mid(y,\widetilde y)\\
 (y,m^*)=(\widetilde y,m^*)=w_2}}
 \Bigg|
 \sum_{\substack{d\\w_0w_1\mid d\\
 (d/w_1,my\widetilde y/w_1^2)=1}}
 \psi_{\Delta_1}(d-d_0)\notag\\[-1mm]
&\hspace{17mm}\times
 \sum_{\substack{n,\widetilde n\\
 y\widetilde n\equiv\widetilde yn\ (d)}}^{\dagger_3}
 C(n;d)C(\widetilde n;d)\psi_N(n)\psi_N(\widetilde n)\notag\\[-1mm]
&\hspace{30mm}\times
 \mathscr K_3(n,\widetilde n;y,\widetilde y,d)
 \Bigg|,
 \label{eq:Sigma3}
\end{align}

where $\dagger_3$ retains

\begin{equation}\label{eq:dagger3}
 (n\widetilde n,w_1r_1q_0u_1v_1v_2)=1,\qquad
 ((n+\ell dr_1)(\widetilde n+\ell dr_1),q_0q_2)=1.
\end{equation}

\begin{lemma}[Supported $q$-van der Corput inequality]
\label{lem:qvdc-supported}
For some admissible $W_1,Y,w_0,w_1,w_2$,

\begin{equation}\label{eq:qvdc-bound}
 \Sigma_2\ll x^{3.5\eta+o(1)}g\Delta\,\Sigma_3^{1/2}.
\end{equation}

Consequently

\begin{equation}\label{eq:Sigma3-target}
 \sup\Sigma_3\ll q_0^2g_0^2N^2x^{-27\eta}
\end{equation}

implies \cref{eq:Sigma2-target}.
\end{lemma}

\begin{proof}
On a fixed $d$, partition the inner sum by
$yn^{-1}\pmod d$.  In \cref{eq:phase-factorization} the $e_d$ factor then
depends only on that residue class and disappears after absolute values.
Retain $d$ squarefree until $w_1=(y,d)$ is fixed.  It follows that $w_1$ is
squarefree and $(w_1,m)=1$.

Only $x^{o(1)}$ values of $w_2$ occur.  Indeed, a nonempty block has
$w_2\le |y|\le x^{C_1}$ for a fixed $C_1$.  If
$m=P_1\cdots P_r$ with $P_1<\cdots<P_r$, replace each $P_j$ by the
$j$-th prime.  This injects the possible divisors of $m^*$ below $x^{C_1}$
into the $p_r$-smooth integers below that bound.  Since $m\le x^{C_*}$ is
squarefree, $r=O(\log x/\log\log x)$ and $p_r=O(\log x)$.  Rankin's
bound with $\sigma=1/\log p_r$ gives

\[
 \#\{w_2:w_2\mid m^*,\ w_2\le x^{C_1}\}
 \le x^{C_1\sigma}\prod_{p\le p_r}(1-p^{-\sigma})^{-1}
 =x^{o(1)}.
\]

Dyadic decomposition in $w_1$ and $y$ and the sum over residue classes cost
at most $x^{2\eta+o(1)}W_1$.  Enlarge the squarefree $d$-sum to the
conditions $w_1\mid d$ and $(d/w_1,w_1)=1$.  The congruence
$yn^{-1}\pmod d$ can then be inverted modulo $d/w_1$.  Replacing
$(n,d)=1$ by the weaker $(n,w_1)=1$ is a nonnegative enlargement.

Cauchy--Schwarz in the residue class and $d$ variables contributes
\[
 x^{\eta/2}\Delta/W_1.
\]
The preceding factor $W_1$ cancels.  The accumulated factor is
$x^{2.5\eta+o(1)}\Delta$.  The two copied
congruences combine to

\[
 y\widetilde n\equiv\widetilde yn\pmod d.
\]

For fixed $y$, the equation $y=h_1v_2^*-h_2v_1^*$ has at most $O(g)$
solutions in $h_1,h_2\sim H^*$: all solutions differ by
$(kv_1^*,kv_2^*)$.  Applying this to $y$ and $\widetilde y$ and then taking
the square root contributes $gx^{\eta/2+o(1)}$.  Finally

\[
 1_{(d/w_1,w_1)=1}
 =\sum_{w_0\mid(d/w_1,w_1)}\mu(w_0)
\]

and the divisor count of the squarefree $w_1$ contributes another
$x^{\eta/2+o(1)}$.  The resulting inner expression is exactly
\cref{eq:Sigma3}, proving \cref{eq:qvdc-bound}.  Taking the square root of
\cref{eq:Sigma3-target} leaves
$q_0g_0Nx^{-13.5\eta}$; multiplication by the prefactor in
\cref{eq:qvdc-bound} gives \cref{eq:Sigma2-target}.
\end{proof}

\subsection{Removal of diagonals and large gcd blocks}

Let

\begin{equation}\label{eq:tT}
 t=(w_2,m),\qquad
 T=\max\{t^{-1},H^{-1}\}
 \frac{x^{\delta+100\eta}H^2N}{g\Delta_1}.
\end{equation}

For fixed $y,\widetilde y$, let $\Sigma_4(y,\widetilde y)$ be the inner
absolute value in \cref{eq:Sigma3}, restricted further by

\begin{equation}\label{eq:retained-gcd}
 \left(
 \frac{y(\widetilde n+Bd)-\widetilde y(n+Bd)}d,m
 \right)\le T.
\end{equation}

\begin{lemma}[Diagonal and large-gcd removal]\label{lem:diagonal-removal}
One has

\begin{align}
 \Sigma_3\ll{}&
 \max\left\{
 \frac{q_0^2x^{\delta+10\eta}H^3}{t},
 \frac{H^4}{t}
 \right\}
 \sup_{y,\widetilde y}\Sigma_4(y,\widetilde y)
 +x^{-47\eta}N^2.
 \label{eq:diagonal-removal}
\end{align}
\end{lemma}

\begin{proof}
By \cref{eq:WY-ranges,eq:UV},

\begin{equation}\label{eq:Y-corrected}
 |Y|\ll\frac{HV}{g}
 \ll\frac{q_0x^{\delta+5\eta}H^2}{g}.
\end{equation}

If a divisor $q$ of the numerator in \cref{eq:retained-gcd} exceeds $T$,
then \cref{eq:Delta1,eq:tT,eq:Y-corrected} give

\begin{equation}\label{eq:Y-over-q}
 \frac{|Y|}{q}
 \ll q_0^{-1}x^{-150\eta}
 \frac{\min\{t,H\}}{H^2}.
\end{equation}

We count the discarded tuples in four disjoint cases.  If
$y=\widetilde y$ and $n=\widetilde n$, their number is

\begin{equation}\label{eq:diag-1}
 O(\Delta H^2N)
 \ll q_0^{-2}x^{-50\eta}N^2.
\end{equation}

If $y=\widetilde y$ but $n\ne\widetilde n$, the congruence with a fixed
$q>T$ and \cref{eq:Y-over-q} give

\begin{equation}\label{eq:diag-2}
 \max\{H^2/t,H\}N^2\frac{|Y|}{q}
 \ll q_0^{-1}x^{-150\eta}N^2.
\end{equation}

If $n=\widetilde n$ and $y\ne\widetilde y$, split according to
$(y-\widetilde y,dq)$.  Divisor expansion bounds the resulting two terms by

\begin{equation}\label{eq:diag-3}
 x^{\delta-144\eta+o(1)}N,
 \qquad
 x^{2\delta-39\eta+o(1)}H^2N.
\end{equation}

The inequalities following from \cref{eq:S52-a,eq:S52-c}, namely
$\gamma>\delta$ and $\gamma>8\omega+4\delta$, together with
$H\ll x^{4\omega+\delta+7\eta}/q_0$, make both terms
$O(x^{-48\eta}N^2)$ after reducing $\eta_0$.

In the fully off-diagonal case, fix $d,\widetilde n,y-\widetilde y$.  The
defining equation leaves $O(q_0\min\{t,H\})$ possible quotient values and
only $x^{o(1)}$ factorizations for each.  Hence this contribution is

\begin{equation}\label{eq:diag-4}
 x^{o(1)}q_0\max\{H^2/t,H\}\Delta N\min\{t,H\}
 \ll q_0^{-1}x^{-50\eta+o(1)}N^2.
\end{equation}

The number of possible gcd divisors is $x^{o(1)}$, so
\cref{eq:diag-1,eq:diag-2,eq:diag-3,eq:diag-4} sum to the error in
\cref{eq:diagonal-removal}.

It remains to count retained phase pairs.  The representation through
$h_1,h_2$ gives

\[
 \#(y,\widetilde y)\ll\max\{H^4/t^2,H^2\},
\]

whereas the interval bound gives
$\#(y,\widetilde y)\ll Y^2/t^2$.  If $t\le x^\delta H$, use the first
estimate; if $t>x^\delta H$, use the second and
\cref{eq:Y-corrected}.  In both cases,

\begin{equation}\label{eq:retained-pairs}
 \#(y,\widetilde y)
 \ll\max\left\{
 \frac{q_0^2x^{\delta+10\eta}H^3}{t},
 \frac{H^4}{t}
 \right\}.
\end{equation}

Multiplying the supremum of $\Sigma_4$ by this count completes the proof.
\end{proof}

\subsection{Relabelling the retained blocks}

Write

\begin{equation}\label{eq:relabel}
 z_1=w_0w_1,\qquad y=w_1\lambda,\qquad
 \widetilde y=w_1\widetilde\lambda,\qquad d=z_1e,
 \qquad \Lambda=Y/w_1.
\end{equation}

An empty or incompatible block is discarded.  In every remaining block,

\begin{gather}
 z_1\ll x^{5\eta}\Delta_1,\qquad (z_1,m)=1,\qquad
 w_1\mid z_1,\qquad
 (\lambda,m^*)=(\widetilde\lambda,m^*)=w_2,
 \label{eq:Z5-arithmetic}\\
 H\ll q_0^{-1}x^{4\omega+\delta+7\eta},\qquad
 x^{-4\eta-\delta}N\ll R\ll x^{-2\eta}N,
 \label{eq:Z5-ranges-a}\\
 x^{1/2-\eta}\ll QR\ll x^{1/2+2\omega+\eta},\qquad
 x^{5\eta}H\ll v_i\ll q_0x^{\delta+5\eta}H,
 \label{eq:Z5-ranges-b}\\
 \frac{RQ^2H}{q_0g\Delta_1}\ll m
 \ll\frac{x^\delta RQ^2H}{g\Delta_1},\qquad
 1\ll|\Lambda|\ll
 \frac{q_0x^{\delta+5\eta}H^2}{w_1g}.
 \label{eq:Z5-ranges-c}
\end{gather}

The apparent asymmetry in the two bounds for $m$ is genuine: the lower
endpoint contains the common factor removed from the first prescribed
divisor, whereas the upper endpoint uses only its ambient dyadic range.

Replace $e$ again by $d$, and set

\begin{equation}\label{eq:relabel-rest}
 \begin{aligned}
 l&=\ell z_1r_1,&
 c_1&=r_1q_0u_1[v_1,v_2],& c_2&=q_0q_2,\\
 A&\leftarrow Aw_1/z_1,&
 B&\leftarrow Bz_1.
 \end{aligned}
\end{equation}

All quotients in \cref{eq:relabel-rest} are units modulo $m$.
The retained sum becomes

\begin{align}
\Sigma_5={}&\Bigg|
 \sum_{(d,m\lambda\widetilde\lambda)=1}
 \psi_{\Delta_1}(z_1d-d_0)
 \sum_{n,\widetilde n}^{\dagger_5}
 C(n;d)C(\widetilde n;d)\psi_N(n)\psi_N(\widetilde n)
 \notag\\[-1mm]
&\hspace{12mm}\times
 e_m\!\left(
 \frac{A\{\lambda(\widetilde n+Bd)-
 \widetilde\lambda(n+Bd)\}}
 {d(n+Bd)(\widetilde n+Bd)}
 \right)\Bigg|,
 \label{eq:Sigma5}
\end{align}

where $\dagger_5$ means

\begin{gather}
 (n\widetilde n,w_1c_1)=1,\qquad
 ((n+ld)(\widetilde n+ld),c_2)=1,
 \label{eq:dagger5-a}\\
 \lambda\widetilde n\equiv\widetilde\lambda n
 \pmod{z_1d/w_1},\qquad
 \left(
 \frac{\lambda(\widetilde n+Bd)-\widetilde\lambda(n+Bd)}d,m
 \right)\le T.
 \label{eq:dagger5-b}
\end{gather}

Equations \cref{eq:relabel,eq:relabel-rest} are reversible on each nonempty
block, so $\Sigma_4=\Sigma_5$.  Combining
\cref{eq:diagonal-removal,eq:Sigma3-target} shows that it is enough to prove

\begin{equation}\label{eq:Sigma5-target}
 \Sigma_5\ll
 \min\left\{\frac{t}{q_0^2x^{\delta+10\eta}H^3},
              \frac{t}{H^4}\right\}
 q_0^2g_0^2N^2x^{-27\eta}.
\end{equation}

\subsection{First elimination: the copied variable}

The congruence in \cref{eq:dagger5-b} is equivalent to

\begin{equation}\label{eq:ntilde-param}
 \widetilde n=
 \frac{\widetilde\lambda n+k(z_1/w_1)d}{\lambda}
\end{equation}

for an integer $k$.  The supports of the two $N$-scale profiles and of the
$\Delta_1$-scale profile imply

\begin{equation}\label{eq:k-range}
 |k|\ll K_0:=\frac{w_1|\Lambda|N}{x^{5\eta}\Delta_1}.
\end{equation}

Put

\begin{equation}\label{eq:S-k}
 S_k=\frac{z_1}{w_1}k+(\lambda-\widetilde\lambda)B.
\end{equation}

For shifted smooth profiles $\psi_N^*,\psi_{\Delta_1}^*$ and residue
classes $d_*,n_*\pmod{q_0}$, define

\begin{align}
\Sigma_6(k):={}&\Bigg|
 \sum_{\substack{d\equiv d_*\ (q_0)\\
 (d,m\lambda\widetilde\lambda)=1}}
 \psi_{\Delta_1}^*(z_1d-d_0)
 \sum_{\substack{n\equiv n_*\ (q_0)\\
 \lambda\mid\widetilde\lambda n+k(z_1/w_1)d}}^{\dagger_6}
 \psi_N^*(n)
 \notag\\[-1mm]
&\hspace{12mm}\times
 e_m\!\left(
 \frac{AS_k}
 {(n+Bd)
  \{\lambda^{-1}(\widetilde\lambda n+k(z_1/w_1)d)+Bd\}}
 \right)\Bigg|,
 \label{eq:Sigma6}
\end{align}

where $\dagger_6$ retains the four coprimality conditions obtained by
substituting \cref{eq:ntilde-param} into \cref{eq:dagger5-a}.

\begin{lemma}[Elimination of $\widetilde n$]\label{lem:remove-ntilde}
If nonnegative weights $\xi_k$ satisfy

\begin{equation}\label{eq:xi-mass}
 \sum_{\substack{|k|\ll K_0\\w_2\mid k\\(S_k,m)\le T}}\xi_k\ll1
\end{equation}

and, uniformly in the profiles and residue classes,

\begin{equation}\label{eq:Sigma6-target}
 \Sigma_6(k)\ll \xi_k
 \min\left\{\frac{t}{q_0^2x^{\delta+10\eta}H^3},
              \frac{t}{H^4}\right\}
 q_0g_0N^2x^{-27\eta},
\end{equation}

then \cref{eq:Sigma5-target} holds.
\end{lemma}

\begin{proof}
Substitution of \cref{eq:ntilde-param} in the phase gives exactly
\cref{eq:S-k,eq:Sigma6}.  For the copied smooth weight, write

\begin{align}
 &\psi_N\!\left(
 \frac{\widetilde\lambda n+k(z_1/w_1)d}{\lambda}
 \right)\notag\\
 &\quad=
 \sum_{j=0}^{J}\frac1{j!}
 \left(\frac{(k/w_1)(z_1d-d_0)}{\lambda N}\right)^j
 \psi^{(j)}\!\left(
 \frac{\widetilde\lambda n+(k/w_1)d_0}{\lambda N}
 \right)+O(x^{-100}),
 \label{eq:Taylor-one}
\end{align}

where $J=\lceil20/\eta\rceil$.  On the support of the sum, the expansion
parameter is at most $x^{-5\eta}$ by \cref{eq:k-range}.  Each summand in
\cref{eq:Taylor-one} is a product of one smooth $N$-profile in $n$ and one
shifted smooth $\Delta_1$-profile in $z_1d-d_0$.  The remainder contributes
$O(x^{-90})$ after trivial summation.

For fixed $k$, the two indicators $C(n;d)$ and
$C(\widetilde n;d)$ are constant after $d,n$ are fixed modulo $q_0$.
There are at most $q_0g_0$ contributing pairs of residue classes.  Moreover,
the divisibility condition in \cref{eq:Sigma6} forces $w_2\mid k$: both
$\lambda$ and $\widetilde\lambda$ are divisible by $w_2$, whereas
$(z_1d/w_1,w_2)=1$.  Summing \cref{eq:Sigma6-target} over $k$ and using
\cref{eq:xi-mass} supplies the missing factor $q_0g_0$, and proves
\cref{eq:Sigma5-target}.  The Taylor remainder is smaller than the target by
the fixed parameter reserve.
\end{proof}

\subsection{Second elimination: the divisibility by
\texorpdfstring{$\lambda$}{lambda}}

Let

\begin{equation}\label{eq:lambda-star}
 \lambda^*=(|\lambda|,|\widetilde\lambda|),\qquad
 a_\lambda=\lambda/\lambda^*,\qquad
 b_\lambda=\widetilde\lambda/\lambda^*.
\end{equation}

By \cref{lem:valuation}, $a_\lambda$ and $b_\lambda$ are units modulo $m$.
We write $\bar a_\lambda,\bar b_\lambda$ for their inverses modulo $m$.
Likewise, $\overline{\widetilde\lambda/w_2}$ denotes the inverse of
$\widetilde\lambda/w_2$ modulo $m$.

If $\lambda^*\nmid(z_1/w_1)k$, then \cref{eq:Sigma6} is empty.  Otherwise
there are residue classes $F_k,G_k$ such that

\begin{equation}\label{eq:n-n1}
 n=a_\lambda n_1+F_kd,\qquad
 \widetilde n=b_\lambda n_1+G_kd.
\end{equation}

They satisfy

\begin{equation}\label{eq:FG-relation}
 \bar b_\lambda(G_k+B)-\bar a_\lambda(F_k+B)
 \equiv
 \bar a_\lambda\,\overline{\widetilde\lambda/w_2}\,
 \frac{S_k}{w_2}\pmod m.
\end{equation}

Define the genuine short scale

\begin{equation}\label{eq:Delta-star}
 \Delta^*=\min\left\{\frac{N}{|\Lambda|x^{5\eta}},\Delta_1\right\}.
\end{equation}

For shifted profiles at scales $\Delta^*/z_1$ and
$|\bar a_\lambda|_{\rm sc}N$ (the latter notation means the real scale
$|\lambda^*/\lambda|N$), let

\begin{align}
\Sigma_7(k):={}&\Bigg|
 \sum_{\substack{d\equiv d_*\ (q_0)\\
 (d,m\lambda\widetilde\lambda)=1}}
 \psi_{\Delta^*/z_1}(d)
 \sum_{n_1}^{\dagger_7}\psi_{|\lambda^*/\lambda|N}(n_1)
 \notag\\[-1mm]
&\quad\times e_m\!\left(
 \frac{\bar a_\lambda\bar b_\lambda AS_k}
 {(n_1+\bar a_\lambda(F_k+B)d)
  (n_1+\bar b_\lambda(G_k+B)d)}
 \right)\Bigg|.
 \label{eq:Sigma7}
\end{align}

Here $\dagger_7$ consists of the congruence
$a_\lambda n_1+F_kd\equiv n_*\pmod{q_0}$ and the four coprimality
conditions obtained from \cref{eq:n-n1}.

\begin{lemma}[Divisibility elimination]\label{lem:divisibility-elimination}
For every retained $k$,

\begin{equation}\label{eq:Sigma6-to-7}
 \Sigma_6(k)
 \ll\frac{\Delta_1}{\Delta^*}
 \sup\bigl(\Sigma_7(k)+O(x^{-90})\bigr),
\end{equation}

where the supremum is over the shifted profiles and the compatible choices
of $F_k,G_k$.
\end{lemma}

\begin{proof}
Dividing the condition
$\lambda\mid\widetilde\lambda n+(z_1/w_1)kd$ by $\lambda^*$ gives one
linear class for $n$ modulo $a_\lambda$ because
$(a_\lambda,b_\lambda)=1$.  Choosing its representative $F_k$ and
substituting gives \cref{eq:n-n1}; the equality defining $G_k$ then gives
\cref{eq:FG-relation} by direct subtraction.

Use a smooth partition of unity to split the original $d$-support into
$O(\Delta_1/\Delta^*)$ intervals of length $\Delta^*/z_1$.  On one such
interval, centered at $x_0$, the only mixed profile is
$\psi_N^*(a_\lambda n_1+F_kd)$.  Its normalized displacement satisfies

\[
 \frac{|F_k|\,|d-x_0|}{N}
 \ll \frac{|\lambda|}{\lambda^*}
       \frac{\Delta^*}{z_1N}
 \ll x^{-5\eta}.
\]

Taylor expansion to order $J$ separates this profile into products of a
shifted $n_1$-profile and a shifted $d$-profile; the remainder is
$O(x^{-90})$ after summation.  The number of partition blocks is the displayed
factor in \cref{eq:Sigma6-to-7}.  No replacement of $\Delta^*$ by
$\Delta_1$ has been made.
\end{proof}

\subsection{Third elimination: all coprimality conditions}

For $q_3\mid m/q_0$, residue classes $d^*,n^*\pmod{q_0q_3}$, shifted
profiles with $N_2\le N$ and $\Delta_2\le\Delta^*$, and units
$A_1,A_2\pmod m$, define

\begin{align}
\Sigma_8(k):={}&\Bigg|
 \sum_{\substack{d\equiv d^*\ (q_0q_3)}}
 \sum_{\substack{n\equiv n^*\ (q_0q_3)}}
 \psi_{\Delta_2/z_1}(d)\psi_{N_2}(n)\notag\\[-1mm]
&\hspace{27mm}\times
 e_m\!\left(
 \frac{w_2A_1L_1}
 {(n+B_1d)(n+(B_1+L_1)d)}
 \right)\Bigg|,
 \label{eq:Sigma8}
\end{align}

where

\begin{equation}\label{eq:L1}
 L_1\equiv A_2\frac{S_k}{w_2}\pmod m,
 \qquad (A_1A_2,m)=1.
\end{equation}

\begin{lemma}[Coprimality elimination]\label{lem:coprimality-elimination}
For every compatible choice in \cref{eq:Sigma7},

\begin{equation}\label{eq:Sigma7-to-8}
 \Sigma_7(k)
 \ll x^{2\eta+o(1)}q_3\sup\Sigma_8(k)+O(x^{-89}).
\end{equation}
\end{lemma}

\begin{proof}
We give the full sequence of finite transformations.  First apply Mobius
inversion to

\[
 (d,(\lambda/w_2)(\widetilde\lambda/w_2))=1,
 \quad(a_\lambda n_1+F_kd,w_1)=1,
 \quad(b_\lambda n_1+G_kd,w_1)=1.
\]

This introduces

\begin{equation}\label{eq:f123}
 f_1\mid(\lambda/w_2)(\widetilde\lambda/w_2),\qquad
 f_2\mid w_1,\qquad f_3\mid w_1,
\end{equation}

and costs $x^{\eta+o(1)}$.  Put

\[
 f_2^*=(f_2,a_\lambda),\qquad
 f_3^*=(f_3,b_\lambda).
\]

Solvability of the two resulting linear congruences first forces a divisor
$f_{23}$ of $d/f_1$.  On common prime factors at which their required
residues disagree, it forces a further squarefree divisor $f_4$.  The
Chinese remainder theorem then gives the exact parametrization

\begin{equation}\label{eq:f-param}
 \begin{aligned}
 d&=f^*d_3,&
 n_1&=w_3n_3+H_kf_4d_3,\\
 f^*&=f_1f_{23}f_4,&
 w_3&=\lcm(f_2/f_2^*,f_3/f_3^*).
 \end{aligned}
\end{equation}

The supports in \cref{eq:f123} and \cref{lem:valuation} give
$(f^*w_3,m)=1$ and $w_3\mid w_1$.  Hence every division in
\cref{eq:f-param} is valid modulo $m$.

The profile in $w_3n_3+H_kf_4d_3$ is the last mixed profile.  On the
$d_3$-support,

\[
 \frac{|H_kf_4|\,\Delta^*}{z_1f^*|\lambda^*/\lambda|N}
 \ll \frac{w_1|\Lambda|\Delta^*}{z_1N}
 \le x^{-5\eta}.
\]

Taylor expansion through order $J$ therefore separates it, with total
remainder $O(x^{-89})$.

Apply Mobius inversion a second time to the five remaining conditions.  It
introduces

\begin{equation}\label{eq:e01234}
 e_0\mid m,\qquad e_1,e_2\mid c_1,\qquad e_3,e_4\mid c_2
\end{equation}

and a further $x^{\eta+o(1)}$ loss.  Let

\[
 g_1=\lcm(q_0,e_0,e_1,e_2,e_3,e_4)=q_0q_3.
\]

Because $q_0,c_1,c_2\mid m$ and $m$ is squarefree,
$q_3\mid m/q_0$.  The congruences for $d_3,n_3$ are either incompatible,
in which case the block is empty, or determine at most

\[
 \frac{g_1}{\lcm(q_0,e_0)}
 \frac{g_1}{\lcm(q_0,e_1,e_2,e_3,e_4)}
 \le\frac{g_1}{q_0}=q_3
\]

pairs of classes modulo $g_1$.  On each pair, collect the two denominator
slopes into $B_1$ and $B_1+L_1$.  The relation
\cref{eq:FG-relation} gives precisely \cref{eq:L1}; multiplying by units
produces $A_1,A_2$.  The resulting sum is \cref{eq:Sigma8}.  Multiplying
the two Mobius costs and the number of residue pairs proves
\cref{eq:Sigma7-to-8}.
\end{proof}

\subsection{Uniformity of the descendant profiles}

\begin{lemma}[Finite-jet closure]\label{lem:finite-jet}
Let $J=\lceil20/\eta\rceil$ and $J_*=3J+2$.  Suppose every input profile
has fixed compact support and satisfies

\begin{equation}\label{eq:input-jets}
 \|\psi^{(r)}\|_1+\|\psi^{(r)}\|_\infty
 \le B_r(\log x)^{A_r},\qquad 0\le r\le J_*.
\end{equation}

Then every terminal profile in \cref{eq:Sigma8} belongs to one common
shifted $C^2$ profile class, independent of all descendant choices.
Moreover

\begin{equation}\label{eq:terminal-polynomial-size}
 m\le x^7,\qquad N_2\le x^7,\qquad \Delta_2/z_1\le x^7
\end{equation}

for sufficiently large $x$.
\end{lemma}

\begin{proof}
Each of the three Taylor expansions differentiates an input profile at most
$J$ times.  Two terminal derivatives therefore require at most
$3J+2=J_*$ input derivatives.  Products, fixed partitions, and affine
changes of variable preserve the compact support and yield uniform $C^2$
seminorms by Leibniz' rule and the chain rule.  Every normalized Taylor
increment is at most $x^{-5\eta}$, so

\[
 x^{-5\eta(J+1)}(\log x)^{O_\eta(1)}=O(x^{-100}).
\]

There are at most $(J+1)^3$ descendants, a number independent of $x$.

For the size assertion, each of
$r_1,q_0,u_1,v_1,v_2,q_2$ is $O(x)$ and

\[
 m=r_1q_0u_1[v_1,v_2]q_2\le x^6.
\]

Also $N_2\le N\le x$ and
$\Delta_2/z_1\le\Delta^*\le\Delta_1\le\Delta\le x$.
The relaxed common exponent seven leaves room for endpoint constants.
\end{proof}

The lemmas in this section form a one-way chain

\begin{equation}\label{eq:descent-chain}
 \Sigma_8\Longrightarrow\Sigma_7\Longrightarrow\Sigma_6
 \Longrightarrow\Sigma_5\Longrightarrow\Sigma_4
 \Longrightarrow\Sigma_3\Longrightarrow\Sigma_2
 \Longrightarrow\Sigma_1^\sharp.
\end{equation}

Every arrow denotes an upper-bound implication, not an asserted equality.
Restrictions are deleted only after an absolute value makes enlargement
legitimate, and no deleted restriction is later restored.

\section{Terminal trace estimate and quantitative closure}
\label{sec:closure}

We now insert the trace estimate from \cref{sec:terminal} into the last sum
of the descent.  This section also retains both branches of $\Delta^*$ and
checks all powers of the common factor $q_0$.

\subsection{The raw terminal bound}

For a retained $k$, put

\begin{equation}\label{eq:Gk}
 G_k=(S_k,m),
 \qquad S_k=\frac{z_1}{w_1}k+(\lambda-\widetilde\lambda)B.
\end{equation}

\begin{lemma}[Terminal estimate and cancellation of $q_3$]
\label{lem:terminal-cancellation}
Uniformly over every descendant in \cref{eq:Sigma8},

\begin{align}
 \Sigma_8(k)
 \ll{}&
 \frac{x^{\eps_{\rm tr}}G_k}{q_0q_3}
 \left(\frac N{\sqrt m}+\sqrt m\right)
 \left(\frac{\Delta^*}{\sqrt m}+\sqrt m\right).
 \label{eq:Sigma8-terminal}
\end{align}

Consequently

\begin{align}
 \Sigma_6(k)
 \ll{}&
 \frac{x^{3\eta}G_k}{q_0}
 \frac{\Delta_1}{\Delta^*}
 \left(\frac N{\sqrt m}+\sqrt m\right)
 \left(\frac{\Delta^*}{\sqrt m}+\sqrt m\right).
 \label{eq:Sigma6-raw}
\end{align}
\end{lemma}

\begin{proof}
Apply \cref{lem:congruence-trace} to \cref{eq:Sigma8} with
$q=q_0q_3$.  This divisor is squarefree by
\cref{lem:prime-support}.  The numerator is $w_2A_1L_1$ and
\cref{eq:L1} gives

\[
 (w_2A_1L_1,m)=(A_1A_2S_k,m)=G_k.
\]

The scale inequalities $N_2\le N$, $\Delta_2/z_1\le\Delta^*$ then give
\cref{eq:Sigma8-terminal}.

Insert this estimate into \cref{eq:Sigma7-to-8}.  Its factor $q_3$ cancels
the denominator $q_3$ exactly; this cancellation occurs before any gcd
averaging.  Then use \cref{eq:Sigma6-to-7}.  Since
$2\eta+\eps_{\rm tr}+o(1)<3\eta$ for sufficiently large $x$, the result is
\cref{eq:Sigma6-raw}.
\end{proof}

\subsection{The gcd average without
\texorpdfstring{$T\ge K$}{T >= K}}

The range \cref{eq:k-range} and the condition $w_2\mid k$ lead to
$k=w_2k_1$ and

\begin{equation}\label{eq:K1}
 |k_1|\le K_1\ll
 \frac{w_1|\Lambda|N}{x^{5\eta}w_2\Delta_1}.
\end{equation}

Because $(z_1/w_1,m)=1$,

\begin{equation}\label{eq:linear-gcd-coeff}
 \left((z_1/w_1)w_2,m\right)=(w_2,m)=t.
\end{equation}

The elementary \cref{lem:gcd-average} therefore gives

\begin{equation}\label{eq:Gk-average}
 \sum_{\substack{|k_1|\le K_1\\G_{w_2k_1}\le T}}
 G_{w_2k_1}
 \le \tau(m)(2tK_1+T).
\end{equation}

No comparison between $K_1$ and $T$ is assumed.  From
\cref{eq:Z5-ranges-c,eq:tT},

\begin{equation}\label{eq:K-over-T}
 \frac{K_1}{T}
 \ll S_x^{O(1)}q_0x^{-100\eta}
 \frac{\min\{t,H\}}{w_2}
 \le S_x^{O(1)}q_0x^{-100\eta}.
\end{equation}

Define

\begin{equation}\label{eq:Lloss}
 \mathcal L=1+K_1/T.
\end{equation}

Since $m\le x^7$, the standard divisor bound gives
$\tau(m)\ll_\eta x^\eta$.  Dividing \cref{eq:Gk-average} by a sufficiently
large constant times $x^\eta tT\mathcal L$ produces nonnegative weights
$\xi_k$ with total mass at most one.  Hence \cref{eq:Sigma6-raw} may be
written

\begin{align}
 \Sigma_6(k)\ll{}&\xi_k\mathcal L
 \frac{x^{4\eta}tT}{q_0}
 \frac{\Delta_1}{\Delta^*}
 \left(\frac N{\sqrt m}+\sqrt m\right)
 \left(\frac{\Delta^*}{\sqrt m}+\sqrt m\right).
 \label{eq:Sigma6-normalized}
\end{align}

The loss $\mathcal L$ appears once.  There is no later Cauchy--Schwarz step
that squares it.

\subsection{The genuine short scale}

The minimum in \cref{eq:Delta-star} cannot be replaced by $\Delta_1$.
Using \cref{eq:Z5-ranges-c} gives

\[
 \frac{N}{|\Lambda|x^{5\eta}}
 \gg\frac{w_1gN}{q_0x^{\delta+10\eta}H^2}.
\]

Together with \cref{eq:Delta1} and the harmless preprocessing losses this
proves, for one absolute $C_0$,

\begin{equation}\label{eq:Delta-star-lower}
 \Delta^*\gg S_x^{-C_0}
 \frac{N}{q_0^2x^{\delta+55\eta}H^2}.
\end{equation}

For the term containing both $\Delta_1$ and $\Delta^*$ we use the exact
identity

\begin{equation}\label{eq:Delta-star-branches}
 \frac1{\Delta_1\Delta^*}
 =\max\left\{
 \frac{|\Lambda|x^{5\eta}}{N\Delta_1},
 \frac1{\Delta_1^2}
 \right\}.
\end{equation}

\subsection{The exponent ledger}

Set

\begin{equation}\label{eq:Bfactor}
 \mathcal B=\min\left\{
 \frac1{q_0^2x^{\delta+10\eta}H^3},
 \frac1{H^4}
 \right\}.
\end{equation}

After substituting $T$ from \cref{eq:tT} into
\cref{eq:Sigma6-normalized}, the bound required in
\cref{eq:Sigma6-target} follows once each of

\begin{equation}\label{eq:three-terminal-terms}
 \frac{N\Delta^*}{m},\qquad N,\qquad m
\end{equation}

is at most

\begin{equation}\label{eq:terminal-RHS}
 \frac{\Delta^*}{\Delta_1}\mathcal B
 \frac{gq_0^2g_0N\Delta_1}{x^{\delta+131\eta}H^2}.
\end{equation}

Indeed, expanding the two completion factors gives

\[
 \left(\frac N{\sqrt m}+\sqrt m\right)
 \left(\frac{\Delta^*}{\sqrt m}+\sqrt m\right)
 =\frac{N\Delta^*}{m}+N+\Delta^*+m,
\]

and $\Delta^*\le N$ absorbs the third summand into the second.

After cancelling \cref{eq:terminal-RHS}, define

\begin{align}
 E_1&=\frac{x^{\delta+131\eta}}{gq_0^2g_0m}
 \max\{q_0^2x^{\delta+10\eta}H^5,H^6\},\notag\\
 E_2&=\frac{x^{\delta+131\eta}}{gq_0^2g_0\Delta^*}
 \max\{q_0^2x^{\delta+10\eta}H^5,H^6\},\notag\\
 E_3&=\frac{mx^{\delta+131\eta}}{gq_0^2g_0N\Delta^*}
 \max\{q_0^2x^{\delta+10\eta}H^5,H^6\}.
 \label{eq:E123}
\end{align}

Thus it suffices to prove $\mathcal LE_j\ll1$ for $j=1,2,3$.

From the upper bound for $\Delta_1$, the lower bound for $m$, and
$RQ^2=x^{-\eta}q_0MH$, one obtains

\begin{equation}\label{eq:m-lower-ledger}
 m\gg\frac{q_0^2x^{54\eta}MH^4}{gN}.
\end{equation}

Using $M\asymp x/N$ and
$H\ll q_0^{-1}x^{4\omega+\delta+7\eta}$, the two branches of $E_1$ are

\begin{equation}\label{eq:E1-ledger}
 E_1\ll
 x^{-1+2\gamma+4\omega+3\delta+94\eta}
 +x^{-1+2\gamma+8\omega+3\delta+91\eta}.
\end{equation}

The lower bound \cref{eq:Delta-star-lower} gives

\begin{equation}\label{eq:E2-ledger}
 E_2\ll S_x^{O(1)}\left(
 x^{28\omega+10\delta-\gamma+245\eta}
 +x^{32\omega+10\delta-\gamma+242\eta}
 \right).
\end{equation}

For $E_3$, the upper bound for $m$ first gives

\begin{equation}\label{eq:m-upper-ledger}
 m\ll\frac{q_0x^{\delta-\eta}MH^2}{g\Delta_1}.
\end{equation}

The two exact alternatives in \cref{eq:Delta-star-branches} satisfy

\begin{align}
 \frac{|\Lambda|x^{5\eta}}{N\Delta_1}
 &\ll\frac{q_0^3x^{2\delta+65\eta}H^4}{N^2},
 \label{eq:branch-Lambda}\\
 \frac1{\Delta_1^2}
 &\ll\frac{q_0^4x^{2\delta+110\eta}H^4}{N^2}.
 \label{eq:branch-Delta}
\end{align}

Substitution in $E_3$ gives four branches:

\begin{align}
 E_3\ll S_x^{O(1)}\bigl(&
 x^{1+44\omega+16\delta-4\gamma+282\eta}
 +x^{1+48\omega+16\delta-4\gamma+279\eta}\notag\\
 &+x^{1+44\omega+16\delta-4\gamma+327\eta}
 +x^{1+48\omega+16\delta-4\gamma+324\eta}
 \bigr).
 \label{eq:E3-ledger}
\end{align}

Before dropping nonpositive factors, the residual powers of $q_0$ in the
eight branches of \cref{eq:E1-ledger,eq:E2-ledger,eq:E3-ledger} are

\begin{equation}\label{eq:q0-ledger}
 -3,-6,-5,-8,-7,-10,-6,-9.
\end{equation}

The second term of $\mathcal L$ in \cref{eq:K-over-T} adds one to each
power of $q_0$ and subtracts $100\eta$ from its power of $x$.  The powers
become

\[
 -2,-5,-4,-7,-6,-9,-5,-8,
\]

and remain nonpositive.  Hence $\mathcal L$ introduces no new parameter
condition.

The $\eta$-free exponents in \cref{eq:E1-ledger,eq:E2-ledger,eq:E3-ledger}
are bounded respectively by

\begin{equation}\label{eq:three-margins}
 -1+8\omega+4\delta+2\gamma,\qquad
 32\omega+10\delta-\gamma,\qquad
 1+48\omega+16\delta-4\gamma.
\end{equation}

They are strictly negative by
\cref{eq:S52-a,eq:S52-b,eq:S52-c}.  The choice
\cref{eq:eta-choice} absorbs every displayed coefficient of $\eta$, the
fixed $S_x^{O(1)}=x^{o(1)}$ losses, and $\eps_{\rm tr}$.  Therefore
$\mathcal LE_j\ll1$ for all three $j$.

\begin{proposition}[Uniform duplicated prescribed-factor estimate]
\label{prop:duplicated-S52}
For every admissible outer localization,

\begin{equation}\label{eq:duplicated-conclusion}
 \Sigma_1^\sharp\ll g_0RQNUV^2x^{-4\eta},
\end{equation}

uniformly in all outer variables and all compatible descendant choices.
\end{proposition}

\begin{proof}
The exponent ledger proves the pointwise estimate
\cref{eq:Sigma6-target} with the normalized weights from
\cref{eq:Sigma6-normalized}.  Apply
\cref{lem:remove-ntilde,lem:diagonal-removal,lem:qvdc-supported,lem:prepared-factor}
in that order.  Uniformity of every trace and Taylor constant follows from
\cref{lem:finite-jet}.  This yields \cref{eq:duplicated-conclusion}.
\end{proof}

\section{Outer Cartesian and dispersion closure}
\label{sec:outer}

This section proves that the local estimate
\cref{prop:duplicated-S52} applies to the individually witnessed modulus
family in \cref{def:S52-family}.  The order in which witnesses are chosen is
part of the argument.

\subsection{Selection and Cartesian enlargement}

Split the modulus range into dyadic blocks $d\asymp D$.  The blocks with

\begin{equation}\label{eq:low-moduli}
 D\le x^{1/2}(\log x)^{-B}
\end{equation}

contribute $O_A(x(\log x)^{-A})$ by the Bombieri--Vinogradov part of the
external dispersion input \ref{input:outer}, after $B=B(A)$ is chosen.
For the remaining blocks, $D\gg x^{1/2-\eta}$ for sufficiently large $x$.

Delete the moduli for which

\begin{equation}\label{eq:exceptional-small-primes}
 \prod_{\substack{p\mid d\\p\le D_0}}p>S_x.
\end{equation}

The standard small-prime-tail estimate in the same external input bounds
their total discrepancy by $O_A(x(\log x)^{-A})$.  For every surviving
modulus, apply \cref{lem:preprocessing}.  Choose the lexicographically first
admissible tuple

\begin{equation}\label{eq:witness-order-first}
 d\longmapsto(q(d),r(d),u(d),d_1(d))
\end{equation}

before any variable is duplicated.  Localizing $q$ and $r$ gives independent
families $\Qfam(Q,R)$ and $\Rfam(Q,R)$ as in \cref{eq:families}, with

\begin{equation}\label{eq:outer-QR}
 x^{\gamma-\delta-3\eta}\ll R\ll x^{\gamma-2\eta},
 \qquad
 x^{1/2-\eta}\ll QR\ll x^{1/2+2\omega}.
\end{equation}

The strict parameter reserve gives

\begin{equation}\label{eq:RQ2}
 RQ^2=\frac{(RQ)^2}{R}
 \ll x^{1+4\omega-\gamma+\delta+3\eta}\ll x.
\end{equation}

Each originally selected pair lies in
$\Qfam(Q,R)\times\Rfam(Q,R)$.  Enlarge to that Cartesian product only after
absolute values have made every summand nonnegative.  Extra cross-pairs and
multiple representations can only enlarge the majorant.  Cross-pairs whose
product is not squarefree are omitted; every selected original pair remains.
No assertion is made that an extra cross-pair inherits either witness.

\subsection{Completion and the common factor}

Let $q^{(1)},q^{(2)}\in\Qfam(Q,R)$ be the two variables produced by
Cauchy--Schwarz and write

\begin{equation}\label{eq:q0-decomposition}
 q_0=(q^{(1)},q^{(2)}),\qquad q^{(i)}=q_0q_i.
\end{equation}

Then $(q_1,q_2)=1$ and

\begin{equation}\label{eq:q0-range}
 1\le q_0\le2Q,\qquad Q/q_0\le q_i\le2Q/q_0.
\end{equation}

Because every member of $\Qfam$ has no prime factor at most $D_0$,

\begin{equation}\label{eq:q0-dichotomy}
 q_0=1\quad\hbox{or}\quad q_0>D_0.
\end{equation}

For the witness $u(q^{(1)})$ chosen in
\cref{eq:witness-order-first}, put

\begin{equation}\label{eq:outer-u1v1}
 u_1=\frac{u(q^{(1)})}{(u(q^{(1)}),q_0)},
 \qquad v_1=q_1/u_1.
\end{equation}

This is the first-factor extraction in \cref{lem:first-factor}; localization
gives \cref{eq:UV}.  If $H<1$, the completed nonzero-frequency sum is empty.
Otherwise localize $h,u_1,v_1$ to scales $H^*,U,V$.

For fixed $q_0,b_1,b_2$ and $0<|\ell|\ll N/R$, let
$\Sigma^\sharp_{H^*,U,V}$ be the completed nonzero-frequency sum before the
second Cauchy--Schwarz inequality.  Its first Cauchy factor is bounded by

\begin{equation}\label{eq:first-Cauchy-factor}
 \Upsilon_1\ll\frac{g_0RQUN}{q_0^2}.
\end{equation}

The second factor is $\Sigma_1^\sharp$, including both copies $v_1,v_2$,
the common variables $r,u_1,q_2$, and all squarefree and coprimality
indicators used in \cref{sec:descent}.  Consequently
\cref{eq:first-Cauchy-factor,eq:duplicated-conclusion} give

\begin{align}
 (\Sigma^\sharp_{H^*,U,V})^2
 &\ll \frac{g_0RQUN}{q_0^2}
       \bigl(g_0RQNUV^2x^{-4\eta}\bigr),\notag\\
 \Sigma^\sharp_{H^*,U,V}
 &\ll\frac{g_0RQN}{q_0}UVx^{-2\eta}
 \ll\frac{g_0RQ^2N}{q_0^2}x^{-2\eta}.
 \label{eq:completed-local}
\end{align}

The last inequality uses $UV\asymp Q/q_0$.  Summing the fixed number of
dyadic and preprocessing localizations changes \cref{eq:completed-local} to

\begin{equation}\label{eq:completed-summed}
 \Sigma^\sharp(Q,R;q_0,b_1,b_2,\ell)
 \ll\frac{g_0RQ^2N}{q_0^2}x^{-2\eta+o(1)}
 \ll\frac{g_0RQ^2N}{q_0^2}x^{-3\eta/2}.
\end{equation}

This is why the local proof keeps a full $x^{-2\eta}$ saving: the outer
localizations consume part, but not all, of it.

The completion is taken at

\[
 H=\frac{x^\eta RQ^2}{q_0M}.
\]

Its modulus is
$k=r[q^{(1)},q^{(2)}]=rq_0q_1q_2\le8RQ^2/q_0$.  Therefore

\[
 \frac{H}{kM^{-1+\eta}}
 \gg (x/M)^\eta\asymp N^\eta\longrightarrow\infty,
\]

so the completion range required in the external dispersion theorem is
valid.  Its discarded-frequency remainder, with a sufficiently large fixed
number of profile derivatives, is

\begin{equation}\label{eq:completion-remainder}
 O_A\left(\frac{MN^2}{R(\log x)^A}\right).
\end{equation}

\subsection{Return to the discrepancy}

Before the final $q_0,\ell$ summation, the nonzero-frequency contribution is
bounded by

\begin{equation}\label{eq:outer-majorant}
 \sum_{1\le|\ell|\ll N/R}\sum_{q_0\le2Q}
 \frac{q_0M}{RQ^2}
 \Sigma^\sharp(Q,R;q_0,b_1,b_2,\ell).
\end{equation}

For every nonzero $\ell$,

\begin{align}
 \sum_{q_0\le2Q}\frac{(q_0,\ell)}{q_0}
 &=\sum_{q_0\le2Q}\frac1{q_0}
   \sum_{c\mid(q_0,\ell)}\varphi(c)\notag\\
 &\le(1+\log(2Q))\sum_{c\mid\ell}\frac{\varphi(c)}c
 \ll\tau(|\ell|)\log(2Q)=x^{o(1)}.
 \label{eq:q0ell-sum}
\end{align}

Insert \cref{eq:completed-summed} into \cref{eq:outer-majorant} and use
\cref{eq:q0ell-sum}.  The result is

\begin{equation}\label{eq:outer-power-saving}
 \ll MNx^{-3\eta/2}
 \sum_{1\le |\ell|\ll N/R}
 \sum_{q_0\le 2Q}\frac{(q_0,\ell)}{q_0}
 \ll\frac{MN^2}{R}x^{-3\eta/2+o(1)}
 \ll\frac{MN^2}{R}x^{-\eta}.
\end{equation}

This is stronger than \cref{eq:completion-remainder} for every prescribed
logarithmic power after enlarging its constant.

The zero frequency with $q_0=1$ is the main term in the dispersion identity.
For $q_0>1$, \cref{eq:q0-dichotomy} gives the standard majorant

\begin{equation}\label{eq:zero-frequency-error}
 MN(\log x)^{O(1)}+
 \frac{MN^2}{RD_0}(\log x)^{O(1)},
\end{equation}

which is absorbed because $R/N\ll x^{-2\eta}$ and $D_0$ dominates every
fixed power of $\log x$.  The diagonal $\ell=0$ is estimated elementarily
using \cref{eq:RQ2} and $R\ll x^{-2\eta}N$.

The initial dispersion identity, input \ref{input:outer}, now converts
\cref{eq:completion-remainder,eq:outer-power-saving,eq:zero-frequency-error}
back to the discrepancy \cref{eq:discrepancy}.  There are only
$(\log x)^{O(1)}$ outer blocks.  Adding the low-modulus and exceptional
small-prime contributions proves \cref{eq:S52-conclusion}, and therefore
\cref{thm:S52}.

All estimates above are uniform in the permitted initial residue $a=a(x)$.
After the initial congruence is fixed, $a$ enters the local argument only
through the unit $A$ in \cref{eq:AB-CRT}; the external complete-sum bounds
and every ensuing absolute-value majorant are uniform in that unit.

The complete order of choices is

\begin{equation}\label{eq:quantifier-order}
 \begin{gathered}
 d\mapsto(q(d),r(d),u(d),d_1(d))
 \mapsto(Q,R,\Qfam,\Rfam),\\
 (Q,R,\Qfam,\Rfam)
 \mapsto(q^{(1)},q^{(2)},q_0)
 \mapsto(u_1,v_1,H^*,U,V).
 \end{gathered}
\end{equation}

In particular, neither witness is allowed to depend on a later Cauchy copy.

\section{The layered variational certificate}
\label{sec:layered}

This section is a finite, computer-assisted theorem.  The analytic
equidistribution theorem proved in the preceding sections is not used in the
evaluation of the integrals.  It enters only when the certified support is
connected to the sieve in \cref{sec:main-proof}.

\subsection{The layered support}

Fix
\begin{equation}\label{eq:finite-parameters}
 k=46,\qquad
 R=\frac{521}{2000},\qquad
 \delta=\frac1{50},\qquad
 \omega=\frac{57}{10000}.
\end{equation}
The endpoint used for the marginal integral is
\begin{equation}\label{eq:outer-radius}
 A_1=\frac14+\omega=\frac{2557}{10000},\qquad
 \varepsilon=R-A_1=\frac3{625},\qquad
 R_o=A_1-\varepsilon=\frac{2509}{10000}.
\end{equation}
Let
\[
\begin{array}{c|cccccccccc}
m&1&2&3&4&5&6&7&8&9&10\\ \hline
B_m&
\frac{157}{1000}&\frac{161}{1000}&\frac{179}{1000}&
\frac{179}{1000}&\frac{188}{1000}&\frac{199}{1000}&
\frac{203}{1000}&\frac{208}{1000}&\frac{208}{1000}&
\frac{211}{1000}.
\end{array}
\]
For definiteness put $B_m=219/1000$ for $m\ge11$.

\begin{definition}[Layered support]\label{def:layered-region}
For $\mathbf t=(t_1,\ldots,t_k)\in[0,1]^k$, set
\[
 L(\mathbf t)=\{i:t_i>\delta\},
 \qquad M(\mathbf t)=|L(\mathbf t)|.
\]
The layered region is
\begin{equation}\label{eq:Omega}
 \Omega=\left\{\mathbf t\in[0,1]^k:
 \sum_{i=1}^k t_i<R,\quad
 \sum_{i\in L(\mathbf t)}t_i\le B_{M(\mathbf t)}
 \right\}.
\end{equation}
\end{definition}

Only the layers $0\le M\le10$ are nonempty.  Indeed,
$10\delta=1/5\le B_{10}$, whereas
$11\delta=11/50>B_{11}$.  The point of \cref{eq:Omega} is that it retains
both the number of coordinates above $\delta$ and their total mass.  A
radial support condition remembers neither datum.

\begin{proposition}[Exact support routing]\label{prop:support-routing}
Every pair of coordinate configurations occurring in the denominator and
the truncated marginal functionals on $\Omega$ satisfies one of the algebraic
support conditions associated with Type~I, Type~IIa, Type~IIb, Type~III,
Bombieri--Vinogradov, or the prescribed-factor condition of
\cref{thm:S52}.  In every prescribed-factor cell, the three strict
parameter inequalities in \cref{thm:S52} hold with positive rational
slack.
\end{proposition}

\begin{proof}
For fixed values of the two large-coordinate counts, all relevant
conditions are rational linear inequalities in the ordered coordinate
masses and the auxiliary partition points.  Of the
$11\cdot12/2=66$ unordered count pairs, the pair $(0,0)$ contains no large
coordinate and is routed directly to the Bombieri--Vinogradov cell.  We
enumerate admissible factor partitions for each of the remaining $65$ pairs
and ask Z3 to test the negation of their union in exact rational
quantifier-free linear arithmetic.  Every returned result is unsatisfiable.
The capacities, all covering partitions, and every strict rational margin
are recorded in the ancillary support certificate.  Thus this is a finite
exhaustive check, rather than a sampled-grid calculation.  The precise
verification interface and its trusted-kernel boundary are described in
\cref{app:certificate}.
\end{proof}

\subsection{The rational test function}

Put $x_i=t_i/R$ and define
\[
 U=\sum_i x_i,\qquad
 U_b=\sum_{i\in L(\mathbf t)}x_i,\qquad
 D=\sum_i x_i^2.
\]
For $1\le m\le10$ let
\begin{equation}\label{eq:Y-m}
 Y_m=\frac{U_b-m\delta/R}{B_m/R-m\delta/R}\in[0,1].
\end{equation}
Let $L_j(y)=P_j(2y-1)$ be the shifted Legendre polynomial.  The frozen
candidate has the form
\begin{equation}\label{eq:test-function}
 F(\mathbf t)=\1_{\Omega}(\mathbf t)
 \prod_{i=1}^{46}q(x_i)
 \left(A_{M(\mathbf t)}(U,Y_{M(\mathbf t)})
       +D\,C_{M(\mathbf t)}(U)\right),
\end{equation}
where the $Y$-dependent terms are omitted in layer $M=0$, and
\begin{equation}\label{eq:A-layer}
 A_m(U,Y)=a_{m,0}(U)+\sum_{\ell=1}^{4}a_{m,\ell}(U)L_\ell(Y).
\end{equation}
Here $q$ is the fixed positive degree-$32$ one-coordinate polynomial whose
$33$ shifted-Chebyshev coefficients, with denominator $2^{60}$, are frozen in
the ancillary profile source identified by hash in
\cref{app:certificate}.  All the $a_{m,\ell}$ and $C_m$ are expanded
in shifted Legendre polynomials in $U$.  Their degree schedules, indexed by
$m=0,\ldots,10$, are
\begin{align}
 \deg a_{m,0}&=(4,4,4,4,4,4,4,3,2,1,1),\notag\\
 \deg a_{m,\ell}\;(\ell\ge1)
 &=(-1,4,4,4,4,3,3,2,1,-1,-1),\notag\\
 \deg C_m&=(3,3,3,3,3,2,2,1,1,-1,-1),
\label{eq:degree-schedules}
\end{align}
where $-1$ means that the mode is absent.  This gives $208$ coefficients.
Every coefficient is an integer divided by $2^{50}$.  The coefficient
ordering and the complete numerator vector are part of the candidate whose
canonical SHA-256 digest is
\begin{equation}\label{eq:candidate-hash}
\begin{gathered}
\mathtt{fbc1d61a9bb8fd7576219001bee8b3c}\\
\mathtt{92728dbfa633d6e710583d8aabd8f98c9}.
\end{gathered}
\end{equation}
The digest binds the $208$ layer coefficients to $k=46$, the degree schedules,
$R_o$, and the support digest.  The polynomial $q$ is bound separately by the
profile-source hash in \cref{app:certificate}.  The floating-point searches
that found these frozen inputs are not used in their certification.

\subsection{The interval calculation}

For a square-integrable function supported on $\Omega$, define
\begin{align}
 I(F)&=\int_{[0,\infty)^k}F(\mathbf t)^2\,\dd\mathbf t,
\label{eq:I-def}\\
 J_i(F)&=
 \int_{\substack{\mathbf t_{\widehat i}\in[0,\infty)^{k-1}\\
                  \sum_{j\ne i}t_j\le R_o}}
 \left(\int_0^\infty F(\mathbf t)\,\dd t_i\right)^2
 \dd\mathbf t_{\widehat i}.
\label{eq:J-def}
\end{align}
By symmetry $J_i(F)$ is independent of $i$; denote it by $J(F)$.
After the change of variables $t_i=Rx_i$, write
\[
 I(F)=R^{46}\overline I,\qquad
 J(F)=R^{47}\overline J.
\]
Notice that \cref{eq:J-def} uses the conservative endpoint $R_o$.
No unverified contribution from its complement is added.

\begin{theorem}[Certified variational inequality]
\label{thm:finite-certificate}
The rational function \cref{eq:test-function} satisfies
\begin{align}
 \overline I&>
 9.0507319290958436555533517458740189572\cdot10^{-110},
\label{eq:I-lower}\\
 \overline J&>
 7.5531779869092232940446089973346969994\cdot10^{-111},
\label{eq:J-lower}\\
 46R\overline J-\overline I&>
 2.4125261747861770030321563214845716954\cdot10^{-114}.
\label{eq:margin-lower}
\end{align}
In particular,
\begin{equation}\label{eq:score-lower}
 \frac{46J(F)}{I(F)}>
 1.0000266555919862183480792294536290427>1.
\end{equation}
\end{theorem}

\begin{proof}
On each layer the integrands are polynomials over a rational polytope.
The exact preprocessing uses rational polynomial arithmetic; the final
integrals are evaluated at $1024$-bit precision with directed outward
rounding in Arb \cite{Arb}.  The $I$ certificate separately encloses all
eleven layer contributions.  The $J$ certificate separately encloses the
small--small, twice-small--large, and large--large channels, and it imports
the certified interval for $I$ by its file hash.  It then evaluates the
signed expression $46R\overline J-\overline I$ directly, rather than
deducing its sign by dividing two rounded decimal approximations.
The resulting interval has the positive lower endpoint displayed in
\cref{eq:margin-lower}.  Candidate and support hashes agree in all three
artifacts.  Full midpoint-radius enclosures and reproduction commands are
given in \cref{app:certificate}.
\end{proof}

\subsection{Removing the polyhedral discontinuity}

The indicator in \cref{eq:test-function} is convenient for exact
integration, whereas the sieve is usually stated for smooth test functions.
The strict margin makes the passage harmless.

\begin{lemma}[Smooth approximation]\label{lem:smooth-approximation}
For every $\epsilon>0$ there is a smooth symmetric function
$F_\epsilon$ compactly supported in the interior of $\Omega$ such that
\[
 |I(F_\epsilon)-I(F)|<\epsilon,\qquad
 |J_i(F_\epsilon)-J_i(F)|<\epsilon
 \quad(1\le i\le46).
\]
Consequently $F_\epsilon$ may be chosen so that
$46J(F_\epsilon)>I(F_\epsilon)$.
\end{lemma}

\begin{proof}
The region $\Omega$ is a finite union of bounded rational polyhedral cells
and its boundary has Lebesgue measure zero.  Thus smooth functions compactly
supported in the cell interiors are dense in $L^2(\Omega)$.  Averaging over
the symmetric group preserves the support and gives symmetric
approximants.  The map
\[
 T_iG(\mathbf t_{\widehat i})=\int_0^\infty G(\mathbf t)\,\dd t_i
\]
has operator norm at most $R^{1/2}$ from $L^2(\Omega)$ to
$L^2(\dd\mathbf t_{\widehat i})$, by Cauchy--Schwarz on each fibre.
Restriction to $\sum_{j\ne i}t_j\le R_o$ cannot increase the norm.
Hence $I(G)=\|G\|_2^2$ and
$J_i(G)=\|\1_{\{\sum_{j\ne i}t_j\le R_o\}}T_iG\|_2^2$
are continuous under $L^2$ approximation.  The positive margin in
\cref{eq:margin-lower} therefore persists.
\end{proof}

Together, \cref{prop:support-routing,thm:finite-certificate,lem:smooth-approximation}
prove \cref{thm:finite-intro}.

\section{The sieve interface and the bound
\texorpdfstring{$H_1\le216$}{H1 <= 216}}
\label{sec:main-proof}

For completeness we isolate the precise special case of the
Maynard--Polymath sieve used here.  This prevents the final implication from
depending on the corresponding proposition in \cite{Stadlmann2026}.  Put
\[
 W=W(x)=\prod_{p\le\log\log\log x}p,
 \qquad B_x=\frac{\varphi(W)}W\log x.
\]
Also set $\theta(n)=\log n$ for $n\in\Pp$ and $\theta(n)=0$ otherwise, and
write
\begin{align*}
 \Delta_\theta(x;a,Q)
 ={}&\sum_{\substack{x+O(1)\le m\le2x+O(1)\\m\equiv a\ (Q)}}\theta(m)\\
 &-\frac1{\varphi(Q)}
 \sum_{\substack{x+O(1)\le m\le2x+O(1)\\(m,Q)=1}}\theta(m).
\end{align*}
If $S\subset[0,\infty)^k$ is compact and downward closed, let
$\mathcal Q_i(S,\epsilon_0;x)$ be the set of squarefree integers
\[
 q=\prod_{j\ne i}[d_j,e_j]
\]
for which the factors $[d_j,e_j]$ are pairwise coprime and coprime to $W$, and
both vectors $(\log_xd_j)_j$ and $(\log_xe_j)_j$ belong to
$(1-\epsilon_0)S$.  We say that $S$ has the \emph{common-residue
prime-distribution property} if, for every $i$, $A>0$, fixed $C>0$, and
integer $a=a(x)$ satisfying $(a,p)=1$ for every prime $p\le x$,
\begin{equation}\label{eq:generated-ed}
 \sum_{q\in\mathcal Q_i(S,\epsilon_0;x)}\tau(q)^C
 \left|\Delta_\theta(x;a,Wq)\right|
 \ll_{A,C,\epsilon_0}\frac{x}{(\log x)^A}.
\end{equation}
The bounded endpoint shifts are uniform.  The harmless divisor weight is
included so that the hypothesis matches the divisor expansion literally;
as usual, it can be recovered from an arbitrarily strong unweighted estimate
by divisor-moment truncation.

\begin{proposition}[Sieve transfer]\label{prop:sieve-transfer}
Let $\mathcal H=\{h_1,\ldots,h_k\}$ be admissible, let $S$ be compact and
downward closed, and let $G\in C_c^\infty(\operatorname{int}S)$ be symmetric.
Assume
\begin{equation}\label{eq:mass-support}
 \sup_{\mathbf t,\mathbf u\in S}
 \sum_{j=1}^k(t_j+u_j)<1
\end{equation}
and suppose that \cref{eq:generated-ed} holds for $S$.  If
\begin{equation}\label{eq:sieve-criterion}
 \sum_{i=1}^kJ_i(G)>I(G),
\end{equation}
then infinitely many translates $n+\mathcal H$ contain at least two primes.
\end{proposition}

\begin{proof}
Choose a residue $b\bmod W$ for which $(b+h_i,W)=1$ for every $i$.
Use the tensor-antiderivative approximation from
\cite[Sections~4--5]{PolymathSieve2014}: after a fixed shrinkage of the
support, write $G$ to arbitrary $C^1$ accuracy as a finite sum of tensor
products $\sum_\ell c_\ell\prod_i g_{\ell,i}$ and put
$f_{\ell,i}(t)=\int_t^\infty g_{\ell,i}(u)\,du$.  With
\[
 \lambda_f(m)=\sum_{d\mid m}\mu(d)f(\log_xd),\qquad
 \nu(n)=
 \left(\sum_\ell c_\ell
 \prod_{i=1}^k\lambda_{f_{\ell,i}}(n+h_i)\right)^2,
\]
    Theorems~3.4 and~3.5 of
\cite{PolymathSieve2014,PolymathSieveErratum2015} give, respectively,
\begin{align}
 \sum_{\substack{x\le n\le2x\\n\equiv b\ (W)}}\nu(n)
 &=(I(G)+o(1))C_x,\label{eq:sieve-mass}\\
 \sum_{\substack{x\le n\le2x\\n\equiv b\ (W)}}
 \nu(n)\theta(n+h_i)
 &=\left(J_i(G)+o(1)\right)C_x\log x
 \qquad(1\le i\le k),
 \label{eq:sieve-prime-mass}
\end{align}
where
\[
 C_x=B_x^{-k}\frac{x}{W}>0.
\]
Condition \cref{eq:mass-support} is exactly the trivial-support hypothesis in
the non-prime asymptotic.  For \cref{eq:sieve-prime-mass}, expansion of the
two divisor sums produces the modulus
$W\prod_{j\ne i}[d_j,e_j]$.  After removing finitely many primes dividing
$\prod_{r<s}(h_r-h_s)$ into $W$, the displayed least common multiples are
pairwise coprime.  The resulting residue class depends on the divisor tuple,
so one must not replace it by a pointwise maximum without proof.

We use the standard rough-residue averaging device.  Put
$P_x=\prod_{p\le x}p$ and let $\mathcal A_i$ be the residue classes modulo
$P_x$ satisfying
\[
 a\equiv b+h_i\pmod W,
 \qquad
 a\equiv h_i-h_j\pmod p
 \quad\text{for one }j\ne i
\]
at each prime $W<p\le x$.  Every $a\in\mathcal A_i$ is coprime to $P_x$.
For a fixed divisor tuple, the generated class modulo
$Wq=W\prod_{j\ne i}[d_j,e_j]$ is represented by
\[
 \frac{|\mathcal A_i|}{(k-1)^{\omega(q)}}
\]
members of $\mathcal A_i$.  Since
$(k-1)^{\omega(q)}\le\tau(q)^{C_k}$ for a fixed $C_k$, its discrepancy is
bounded by
\[
 \frac{\tau(q)^{C_k}}{|\mathcal A_i|}
 \sum_{a\in\mathcal A_i}|\Delta_\theta(x;a,Wq)|.
\]
After summing the divisor tuples, their multiplicity is another fixed power
of $\tau(q)$.  Averaging \cref{eq:generated-ed} over
$a\in\mathcal A_i$ therefore bounds the full error by
$O_A(x(\log x)^{-A})$.  This is the residue-class transfer used in
\cite[Section~2.4]{Stadlmann2026}.  The main term is the one-coordinate
deletion integral $J_i(G)$.

By \cref{eq:sieve-criterion,eq:sieve-mass,eq:sieve-prime-mass}, for all
sufficiently large $x$ one has
\[
 \sum_{\substack{x\le n\le2x\\n\equiv b\ (W)}}\nu(n)
 \left(\sum_{i=1}^k\theta(n+h_i)-\log(3x)\right)>0.
\]
Since $\nu(n)\ge0$ and $n+h_i<3x$ for large $x$, some $n$ in the
progression has at least two prime translates.
The same argument on every sufficiently large dyadic interval gives
infinitely many such translates.
\end{proof}

Proposition~\ref{prop:support-routing} gives the exact support routing.  To
make the analytic dependencies auditable, \cref{tab:support-inputs} records
the estimate used for each named cell.  Lemma numbers in the non-BV rows
refer to version~1 of \cite{Stadlmann2026}; those lemmas are themselves
relaxed-modulus variants of the published results shown in the right column.
\begin{table}[ht]
\small
\begin{tabular}{@{}lll@{}}
\toprule
cell & direct input & antecedent or replacement\\
\midrule
BV & Bombieri--Vinogradov & \cite{PolymathSieve2014}\\
Type I & \cite[Lemma~5]{Stadlmann2026} & \cite[Lemma~5]{BakerIrving2017}\\
Type IIa & \cite[Lemma~3]{Stadlmann2026} &
  \cite[Theorem~2.8(iv)]{PolymathEDZ2014}\\
Type IIb & \cite[Lemma~4]{Stadlmann2026} &
  \cite[Theorem~2.8(ii)]{PolymathEDZ2014}\\
Type IIc & \cref{thm:S52} & replaces
  \cite[Lemma~6]{Stadlmann2026}\\
Type III & \cite[Lemma~7]{Stadlmann2026} &
  \cite[Theorem~2.8(v)]{PolymathEDZ2014}\\
\bottomrule
\end{tabular}
\caption{Cell-by-cell analytic inputs for the exact support router.}
\label{tab:support-inputs}
\end{table}
The support certificate verifies the hypotheses of these six rows for every
pair of cells.  The non-Type-IIc rows are the established relaxed-modulus
forms cited in the table, while the Type-IIc row is \cref{thm:S52}, proved in
\cref{sec:terminal,sec:interfaces,sec:descent,sec:closure,sec:outer}.
The standard finite convolution decomposition used in
\cite[Sections~3--4]{Stadlmann2026}, followed cell by cell through the exact
support router, therefore gives \cref{eq:generated-ed}.  The harmless factor
$W=x^{o(1)}$ is absorbed by the small-prime preprocessing.  Finally, the
inequality $2R<1$ gives \cref{eq:mass-support}.  Hence the smooth approximant
furnished by \cref{lem:smooth-approximation} has the common-residue
prime-distribution property required by \cref{prop:sieve-transfer}.

\begin{lemma}[A diameter-\(216\) admissible tuple]\label{lem:tuple}
The set
\begin{align*}
\mathcal H=\{&
0,4,6,16,18,28,30,34,40,48,54,58,60,64,70,76,78,84,88,96,\\
&100,106,114,118,120,126,130,138,144,148,154,160,166,168,\\
&174,180,184,186,190,196,198,204,208,210,214,216\}
\end{align*}
is an admissible \(46\)-tuple of diameter \(216\).
\end{lemma}

\begin{proof}
The entries are distinct and lie in $[0,216]$.  For every prime
$p\le43$, direct reduction leaves at least one residue class unoccupied;
the complete missing-residue table appears in \cref{app:certificate}.
For $p>46$, a set of $46$ integers cannot occupy all $p$ residue classes.
Hence no prime sees every residue class and $\mathcal H$ is admissible.
\end{proof}

\begin{proof}[Proof of \cref{thm:main}]
By \cref{thm:finite-certificate,lem:smooth-approximation}, choose a smooth
symmetric admissible test function $G$ satisfying
\[
 \sum_{i=1}^{46}J_i(G)=46J(G)>I(G).
\]
The support inputs in \cref{tab:support-inputs}, including the proved
\cref{thm:S52}, verify the distribution property of
\cref{prop:sieve-transfer}.  Apply that proposition to the
tuple in \cref{lem:tuple}.  Infinitely many intervals of length $216$
therefore contain two primes.  Replacing those two primes by two consecutive
primes lying between them can only decrease their distance, so
\[
 H_1=\liminf_{n\to\infty}(p_{n+1}-p_n)\le216.
\]
\end{proof}

\begin{remark}[Logical status]\label{rem:verification-status}
No conjectural equidistribution hypothesis is used in the theorem above.
The Type-IIc input is proved as \cref{thm:S52}; the remaining analytic rows
invoke the established results cited in \cref{tab:support-inputs}.  The
$k=46$ variational inequality and tuple admissibility are finite theorems
supported by exact and outward-rounded verification artifacts.  Thus the
conclusion $H_1\le216$ is unconditional in the usual mathematical sense.
\end{remark}


\section*{Acknowledgements and computational disclosure}
This work developed through an extended human--AI collaboration.  The author
formulated the prescribed-factor research program, identified the need to
retain factor-address information, and obtained the initial diameter-$226$
construction.  The antecedent prescribed-factor framework and its proposed
diameter-$240$ application are due to Stadlmann; the two analytic inheritance
issues in that proposal are identified and repaired in this paper.
OpenAI Codex proposed the layered address-state variational ansatz that reduced
the finite candidate diameter from $226$ to $216$, helped design and implement
the exact-support and interval-arithmetic verification pipeline,
assisted in diagnosing and repairing the two source-level inheritance issues
treated in \cref{sec:diagnosis}, and supported the organization and preparation of the
manuscript.  The author reviewed the resulting arguments and artifacts and
accepts full responsibility for every mathematical claim, citation, and
conclusion in the paper.  The finite certificates and source code are supplied
so that the computer-assisted component can be checked independently.

\appendix
\section{Corrections, attribution, and verification}
\label{sec:status}

The logical dependency graph of the argument is

\begin{equation}\label{eq:dependency-DAG}
 \begin{gathered}
 \text{published dispersion and trace-function inputs},\\
 \text{the prescribed-factor and support lemmas of \cref{sec:interfaces}},\\
 \text{the expanded local descent of \cref{sec:descent}},\\
 \text{the branch ledger of \cref{sec:closure}}\\[1mm]
 \Downarrow\\[1mm]
 \cref{prop:duplicated-S52}
 \quad\Longrightarrow\quad
 \cref{thm:S52}.
 \end{gathered}
\end{equation}

The final equidistribution conclusion is never used to prove a local
estimate.  In particular, the outer Cartesian enlargement in
\cref{sec:outer} consumes \cref{prop:duplicated-S52} only after the latter
has been proved.  Thus \cref{eq:dependency-DAG} is acyclic.

The argument uses three kinds of material, which should not be conflated.

\begin{enumerate}[label=\textup{(\arabic*)}]
\item The Deligne-type complete-sum estimates and the initial dispersion
identity are cited established theorems, not reproved here.

\item The order of the $q$-van der Corput transformations originates in
Stadlmann's argument.  The exact prescribed-factor family with its
$52\eta$ window was recorded in
\cite{Stadlmann2026}.  We retain both attributions even though the middle
calculation has been rewritten rather than incorporated by reference.

\item The corrected prescribed-factor family, reversible small-prime
preprocessing, prime-support transport, full $q_0$ ledger, genuine
$\Delta^*$ bifurcation, and gcd average without $T\ge K$ are the modified
parts proved in this document.
\end{enumerate}

For reproducibility, the following finite checks accompany the manuscript:

\begin{center}
\begin{tabularx}{\textwidth}{@{}>{\raggedright\arraybackslash}X
  >{\raggedright\arraybackslash}p{0.38\textwidth}@{}}
\toprule
check & purpose\\
\midrule
\path{supplement/scripts/check_bounded_gap_s52_support_transport.py}
 & squarefree support and CRT units\\
\path{supplement/scripts/check_bounded_gap_s52_reversible_preprocessing.py}
 & witness transport under preprocessing\\
\path{supplement/scripts/check_bounded_gap_s52_q0_transport.py}
 & full $q_0$ loss transport\\
\path{supplement/scripts/check_bounded_gap_s52_gcd_repair.py}
 & finite audit of the gcd-average repair\\
\path{supplement/scripts/check_bounded_gap_s52_exponent_ledger.py}
 & rational exponent and $q_0$ bookkeeping\\
\path{supplement/scripts/check_bounded_gap_s52_outer_closure.py}
 & outer Cauchy and dispersion monomials\\
\bottomrule
\end{tabularx}
\end{center}

These programs are regression checks only; none replaces an analytic
estimate or a proof above.

Relative to the established inputs \ref{input:trace}--\ref{input:qvdc}, the
displayed argument proves \cref{thm:S52}; no conjectural estimate is used.
For clarity, the proof has been organized around the following four
verification checkpoints:

\begin{enumerate}[label=\textup{(V\arabic*)}]
\item verify that the stated form of the squarefree trace estimate follows
from the cited sheaf results under both denominator conventions;
\item compare every support deletion in \cref{sec:descent} against the
frozen source calculation, including empty subunit ranges;
\item independently recompute the eight exponent branches and the
$\mathcal L$ correction;
\item verify the exact hypotheses of the external dispersion reduction for
the Cartesian families in \cref{sec:outer}.
\end{enumerate}

These checkpoints record where each cited input enters; they are not
additional hypotheses of \cref{thm:S52}, and the finite programs above are
not substitutes for the analytic argument.
This paper establishes the prescribed-factor Type-I estimate itself and then supplies
the separate sieve optimization and exact support check required for its
bounded-gap application.  Those components are combined in
\cref{sec:main-proof} to prove \cref{thm:main}.

\section{Finite certificates and reproduction}
\label{app:certificate}

\subsection{Frozen artifacts}

The finite verification uses the five output artifacts and the separately
hashed one-coordinate profile source in \cref{tab:artifact-manifest}.  The
first column is the SHA-256 digest of the file bytes.  The canonical
candidate digest in \cref{eq:candidate-hash} is different: it hashes only
the mathematical binding object and is stable under harmless JSON
formatting changes.

\begin{table}[ht]
\small
\begin{tabular}{@{}p{\textwidth}@{}}
\toprule
\textbf{Frozen file and file-level SHA-256 digest}\\
\midrule
\path{bounded_gap_216_address_state_rational_candidate.json}\\[-2pt]
\quad\texttt{c99e6bfc63fd530250ab3d439ff688c9}\allowbreak
\texttt{e9c4eb9243bf04c37d6b895f15885fb2}\\[3pt]
\path{bounded_gap_216_address_state_support_exact.json}\\[-2pt]
\quad\texttt{18c4c52703f3415d6cb59e02259cb832}\allowbreak
\texttt{f700a0210fbc5d46dd4952417964ed0c}\\[3pt]
\path{bounded_gap_216_address_state_i_arb.json}\\[-2pt]
\quad\texttt{7a95817358ecace2879270aee77926795}\allowbreak
\texttt{bcaeeb80ab38f7bf7f545bea0d9f405}\\[3pt]
\path{bounded_gap_216_address_state_j_arb.json}\\[-2pt]
\quad\texttt{fa7a7124992aba6014483ed10ec8908b}\allowbreak
\texttt{26ad598579933bf6f1dbd29049029556}\\[3pt]
\path{bounded_gap_216_single_assumption_audit.json}\\[-2pt]
\quad\texttt{4181126c1c84f366f03ea43894fba53a}\allowbreak
\texttt{fce2c0aeff4828cbe44ff844104d83b5}\\[3pt]
\path{build_bounded_gap_236_rational_candidate.py}\\[-2pt]
\quad\texttt{e3d44a26ec147c5eca0b04491ee044a9}\allowbreak
\texttt{ea41d6b20099364bf2e98a340d2b3bcd}\\
\bottomrule
\end{tabular}
\caption{Finite-certificate and profile-source manifest.}
\label{tab:artifact-manifest}
\end{table}

The ancillary package also contains the machine-readable file
\path{MANIFEST.sha256}, which covers \path{requirements.txt}, the ancillary
README, every supplied Python source, and every frozen JSON output.  Its
SHA-256 digest is
\[
 \begin{gathered}
 \mathtt{ce39da06a82645414027c5a0c701982c}\\
 \mathtt{eefb5898c171a4f3443806a8cdad9f3a}.
 \end{gathered}
\]
Generated virtual environments, bytecode, and cache directories are
deliberately excluded from the package.

The support calculation is performed in exact rational arithmetic.  Its
current proof kernel trusts the QF-LRA unsatisfiability answers returned by
Z3 \cite{Z3}; the JSON records the complete rational input and the covering
partitions, but not a separately replayable solver proof object.  The
integral calculation uses FLINT rational polynomials \cite{FLINT} followed by
Arb balls with directed outward rounding \cite{Arb}.  No floating-point
optimizer output enters the proof.

\subsection{Reproduction commands}

From the ancillary package root, create an isolated environment and install
the versions recorded in \path{requirements.txt}:
\begin{lstlisting}[language=bash,basicstyle=\ttfamily\footnotesize,
breaklines=true]
shasum -a 256 --check MANIFEST.sha256
python3 -m venv .venv
.venv/bin/python -m pip install -r requirements.txt
\end{lstlisting}
The frozen finite candidate is then independently replayed and certified by
\begin{lstlisting}[language=bash,basicstyle=\ttfamily\footnotesize,
breaklines=true]
PYTHONPATH=scripts .venv/bin/python \
  scripts/check_bounded_gap_226_layered_support.py \
  --k 46 \
  --candidate-json \
  outputs/bounded_gap_216_address_state_rational_candidate.json \
  --output \
  outputs/bounded_gap_216_address_state_support_exact.json

PYTHONPATH=scripts .venv/bin/python \
  scripts/certify_bounded_gap_216_address_state_i_arb.py \
  --precision 1024

PYTHONPATH=scripts .venv/bin/python \
  scripts/certify_bounded_gap_216_address_state_j_arb.py \
  --precision 1024

PYTHONPATH=scripts .venv/bin/python \
  scripts/audit_bounded_gap_216_single_assumption.py
\end{lstlisting}
The filename of the generic support checker retains the earlier
\texttt{226} development label; the explicit \texttt{--k 46} and candidate
arguments bind this run to the present certificate.
The floating-point optimization history that discovered the frozen candidate
is not part of the ancillary package.  The absolute path in its provenance
field is descriptive metadata only and is excluded from the canonical
mathematical digest.  Reproduction here means replaying the exact support and
interval certificates from the frozen rational input, not reproducing the
heuristic search that found it.

The analytic regression checks are
\begin{lstlisting}[language=bash,basicstyle=\ttfamily\footnotesize,
breaklines=true]
.venv/bin/python scripts/check_bounded_gap_s52_exponent_ledger.py
.venv/bin/python scripts/check_bounded_gap_s52_q0_transport.py
.venv/bin/python scripts/check_bounded_gap_s52_gcd_repair.py
.venv/bin/python scripts/check_bounded_gap_s52_reversible_preprocessing.py
.venv/bin/python scripts/check_bounded_gap_s52_support_transport.py
.venv/bin/python scripts/check_bounded_gap_s52_loss_budget.py
.venv/bin/python scripts/check_bounded_gap_s52_finite_jet_closure.py
.venv/bin/python scripts/check_bounded_gap_s52_outer_closure.py
\end{lstlisting}
These scripts verify finite algebraic identities and artifact bindings.
They do not replace the analytic arguments in
\cref{sec:terminal,sec:interfaces,sec:descent,sec:closure,sec:outer}.

\subsection{Admissibility table}

For the tuple in \cref{lem:tuple}, the following residue classes are
unoccupied.  One missing class for each prime would suffice; all classes
reported by the audit are retained here.
\[
\begin{array}{c|l@{\qquad}c|l}
p&\text{missing residues}&p&\text{missing residues}\\
\hline
2&1&3&2\\
5&2&7&3\\
11&2&13&7\\
17&5&19&17\\
23&21&29&8,17\\
31&1,15&37&8,24,35\\
41&1,8,12,27,33,39&
43&3,7,13,22,23,29
\end{array}
\]
For every prime \(p>46\), cardinality alone supplies an unoccupied class.

\subsection{Source provenance}

For auditing the extension against the source arguments, the frozen source
snapshots used in development had digests
\[
\begin{array}{c}
\text{Polymath equidistribution source}\\
\mathtt{fdffe1dfb7b820d8a45ecc0e07e2f7e}\\
\mathtt{17404e6e10b63db110c2d44afe42013ea}\\[3pt]
\text{Stadlmann smooth-modulus source}\\
\mathtt{60c0440f33d9cbf504470716491fb4d4}\\
\mathtt{5b45b26d9a960c8e34ff2af500837a30}\\[3pt]
\text{Stadlmann 2026 version-1 source}\\
\mathtt{c0d5d2317c77f4de7eacdef6e1d4b1e}\\
\mathtt{b6433e6240b5c09273b3d4eee99e6c3ba}
\end{array}
\]
The first two rows are established inputs.  The third is used to identify and
compare the proposed relaxed-modulus estimates and their support conditions.
The present paper proves the corrected prescribed-factor theorem
\cref{thm:S52}; the source snapshot is not used as authority for that
conclusion.


\bibliographystyle{amsplain}
\bibliography{references}

\end{document}
